\chapter{鏈結串列}
\label{ch:lists}

% 本章的圖共用的樣式：一個節點畫成「資料｜next」兩格。
\tikzset{
    lkey/.style={draw=black!60, minimum width=0.85cm, minimum height=0.65cm,
        font=\ttfamily\small, inner sep=1pt},
    lptr/.style={draw=black!60, minimum width=0.40cm, minimum height=0.65cm, inner sep=0pt},
    lhl/.style={lkey, draw=memStroke, thick, fill=memFill},
    lhd/.style={lkey, fill=black!8},
    lname/.style={font=\small\ttfamily},
    lnk/.style={-{Stealth[length=1.8mm]}, thick, draw=black!70},
    lnew/.style={lnk, draw=memStroke},
    lold/.style={lnk, densely dashed, draw=black!30},
    stepno/.style={circle, draw=optStroke, fill=optFill, inner sep=0.5pt,
        font=\scriptsize\bfseries, text=optStroke},
}
% \lnode[樣式]{名稱}{座標}{內容}：資料格叫 <名稱>，next 格叫 <名稱>p。
\newcommand{\lnode}[4][lkey]{%
    \node[#1] (#2) at (#3) {#4};
    \node[lptr, anchor=west] (#2p) at (#2.east) {};
}
% \lnil{名稱}：next 是 NIL，在 next 格畫一條斜線。
\newcommand{\lnil}[1]{\draw[black!60] (#1p.south west) -- (#1p.north east);}
% \llink[樣式]{名稱}{路徑}：從 <名稱> 的 next 格畫出指標，例如 \llink{a}{-- (b)}。
\newcommand{\llink}[3][lnk]{%
    \fill[black!70] (#2p.center) circle (1.1pt);
    \draw[#1] (#2p.center) #3;
}

社團的迎新活動開放報名，總召把報名的同學依學號由小到大寫在一張紙上。 每天都有
人新報名，也有人取消。 新報名的人要插進中間某個位置，名單如果寫在一整排連續的格子
裡，插入一個人就得把後面所有的人往後抄一格。 換一種記法：每個人旁邊只註明「下一位
是誰」，插入時只要改一兩個「下一位」。 名單有 $n$ 個人時，兩種記法插入一個人各要
改幾個地方？ 第二種記法有沒有什麼事變得比較慢？ 往下讀之前，先想一想。

第二種記法就是\textbf{鏈結串列}（linked list）。 這一章先講單向串列的搜尋、插入和
刪除，再講幾種變形：帶頭節點的串列、環狀串列、雙向串列。 接著用串列表示多項式和稀
疏矩陣，說明怎麼自己管理一池節點，再整理反轉、合併、找中間這三個最常考的操作，最後
用快慢指標判斷串列有沒有環、找出環從哪裡開始。 每個操作都會給虛擬碼、C 程式、手算
追蹤和成本。

這一章用到《C 程式設計基礎》第~14~章的結構指標（\texttt{p->next}）、第~19~章的
\texttt{malloc}、\texttt{free} 和串列入門。 那一章已經寫過頭端插入、指標的指標和
刪除；這一章從資料結構的角度重新整理，重點放在成本、正確性和考卷上的追蹤，用到時
會先在框裡複習。

\section{單向鏈結串列}
\label{sec:ll-singly}

\subsection{節點與串列}

鏈結串列由\textbf{節點}（node）組成。 單向串列（singly linked list）的每個節點有
兩個欄位：$\mathit{key}$ 存資料，$\mathit{next}$ 指向下一個節點；最後一個節點的
$\mathit{next}$ 是 \NIL。 串列本身只記一個指標 $L.\mathit{head}$，指向第一個節點；
空串列的 $L.\mathit{head} = \NIL$。 實際的節點常常還有別的資料欄位（姓名、成績），
虛擬碼只寫 $\mathit{key}$ 和指標，其他欄位跟著節點一起搬。

\begin{center}
\begin{tikzpicture}
\node[lname] (hd) at (0,0) {head};
\lnode{a}{1.4,0}{115}
\lnode{b}{3.4,0}{160}
\lnode{c}{5.4,0}{207}
\lnode{d}{7.4,0}{342}
\draw[lnk] (hd) -- (a);
\llink{a}{-- (b)}
\llink{b}{-- (c)}
\llink{c}{-- (d)}
\lnil{d}
\end{tikzpicture}
\diagcaption{有四個節點的單向串列。 每個節點左格是 $\mathit{key}$，右格是
$\mathit{next}$；斜線表示 \NIL。}
\label{fig:ll-singly}
\end{center}

節點在記憶體裡不必相鄰，每個節點各自用 \texttt{malloc} 要來，彼此只靠
$\mathit{next}$ 串起來。 這帶來兩個後果。 好處是：只要手上有某個節點 $p$，在它後面
接一個新節點只要改兩個指標，和串列多長無關。 壞處是：想拿第 $k$ 個節點，只能從
$\mathit{head}$ 開始一個一個往下走，沒有辦法像陣列那樣用位址公式直接算出來（第~4~
章）。 下表把兩者並排。

\begin{center}
\begin{tabular}{lll}
\toprule
操作（$n$ 個元素） & 陣列（依序存放） & 單向串列（只記 $\mathit{head}$） \\
\midrule
取第 $k$ 個 & $\Theta(1)$ & $\Theta(k)$ \\
找某個 $\mathit{key}$（沒有排序） & $\Theta(n)$ & $\Theta(n)$ \\
找某個 $\mathit{key}$（已經排序） & $\Theta(\lg n)$，二分搜尋 & $\Theta(n)$ \\
在最前面插入 & $\Theta(n)$，全部往後挪 & $\Theta(1)$ \\
已知節點 $p$，在它後面插入或刪除 & $\Theta(n)$ & $\Theta(1)$ \\
額外記憶體 & 沒有（但要先開好容量） & 每個節點一個指標 \\
\bottomrule
\end{tabular}
\end{center}

排好的串列不能二分搜尋：二分搜尋要直接跳到中間那一格，串列走到中間就要
$\Theta(n)$ 步。 所以「串列比陣列好」或「陣列比串列好」都不對，要看哪些操作比較常
做。 第~6~章的堆疊與佇列、第~12~章雜湊的鏈結法、第~21~章圖的鄰接串列，都是挑串列
擅長的操作來用。

\subsection{搜尋、插入、刪除}

\textbf{搜尋}從 $\mathit{head}$ 往下走，遇到 $\mathit{key}$ 相同的節點就停。

\begin{pseudo}[label=alg:ll-search]{在串列裡找 $k$}
\Function{List-Search}{$L, k$}
    \State $x \gets L.\mathit{head}$
    \While{$x \ne \NIL$ \AND $x.\mathit{key} \ne k$}
        \State $x \gets x.\mathit{next}$
    \EndWhile
    \State \Return $x$ \Comment{找不到時是 \NIL}
\EndFunction
\end{pseudo}

最壞情況（$k$ 不在串列裡）要走過全部 $n$ 個節點，$\Theta(n)$。 第 3 行的
\textbf{and} 先檢查 $x \ne \NIL$，$x$ 是 \NIL 時不會去讀 $x.\mathit{key}$；C 的
\texttt{\&\&} 也是這樣，兩個條件的順序不能對調。

\textbf{插入}要先說好插在哪裡。 單向串列的節點只知道下一個是誰，不知道上一個是誰，
所以插入和刪除都以「前一個節點」$p$ 為準：把新節點 $x$ 接在 $p$ 後面，或拿掉 $p$
後面那一個。 $p = \NIL$ 代表「在最前面」。

\begin{pseudo}[label=alg:ll-insert]{把節點 $x$ 接在 $p$ 後面}
\Procedure{List-Insert-After}{$L, p, x$}
    \If{$p = \NIL$} \Comment{接在最前面}
        \State $x.\mathit{next} \gets L.\mathit{head}$
        \State $L.\mathit{head} \gets x$
    \Else
        \State $x.\mathit{next} \gets p.\mathit{next}$
        \State $p.\mathit{next} \gets x$
    \EndIf
\EndProcedure
\end{pseudo}

\begin{pseudo}[label=alg:ll-delete]{拿掉 $p$ 後面的節點}
\Function{List-Delete-After}{$L, p$}
    \If{$p = \NIL$} \Comment{拿掉第一個}
        \State $x \gets L.\mathit{head}$
        \State $L.\mathit{head} \gets x.\mathit{next}$
    \Else
        \State $x \gets p.\mathit{next}$
        \State $p.\mathit{next} \gets x.\mathit{next}$
    \EndIf
    \State \Return $x$ \Comment{由呼叫的人決定要不要釋放}
\EndFunction
\end{pseudo}

兩個程序都是 $\Theta(1)$：不管串列多長，都只改兩個指標。 前提是被拿掉的節點存在
（$p = \NIL$ 時串列不是空的，否則 $p.\mathit{next} \ne \NIL$）。 實際上要刪除某個
$\mathit{key}$ 時，還得先花 $\Theta(n)$ 找到它的前一個節點，所以「串列刪除是
$\Theta(1)$」只在已經拿到前一個節點時成立。

插入的兩行順序不能對調。 下圖把 $\mathit{key} = 180$ 的節點 $x$ 接在 $160$ 後面。
如果先做 $p.\mathit{next} \gets x$，$207$ 那個節點就再也沒有人指著，接下來的
$x.\mathit{next} \gets p.\mathit{next}$ 讀到的是 $x$ 自己，$x$ 指向自己，後半段整條
遺失。

\begin{center}
\begin{tikzpicture}
% 插入前
\node[anchor=east, font=\small] at (-0.2,0) {插入前};
\lnode{a}{0.6,0}{115}
\lnode[lhl]{b}{2.6,0}{160}
\lnode{c}{4.6,0}{207}
\lnode{d}{6.6,0}{342}
\llink{a}{-- (b)}
\llink{b}{-- (c)}
\llink{c}{-- (d)}
\lnil{d}
\node[font=\scriptsize, text=memStroke] at (2.8,0.6) {$p$};
% 第一步
\node[anchor=east, font=\small] at (-0.2,-1.9) {第 6 行};
\lnode{a2}{0.6,-1.6}{115}
\lnode[lhl]{b2}{2.6,-1.6}{160}
\lnode{c2}{4.6,-1.6}{207}
\lnode{d2}{6.6,-1.6}{342}
\lnode{x2}{3.6,-2.5}{180}
\llink{a2}{-- (b2)}
\llink{b2}{-- (c2)}
\llink{c2}{-- (d2)}
\lnil{d2}
\llink[lnew]{x2}{-- (c2.south)}
\node[font=\scriptsize, text=black!60] at (2.6,-2.5) {$x$};
% 第二步
\node[anchor=east, font=\small] at (-0.2,-4.1) {第 7 行};
\lnode{a3}{0.6,-3.8}{115}
\lnode[lhl]{b3}{2.6,-3.8}{160}
\lnode{c3}{4.6,-3.8}{207}
\lnode{d3}{6.6,-3.8}{342}
\lnode{x3}{3.6,-4.7}{180}
\llink{a3}{-- (b3)}
\fill[black!70] (b3p.center) circle (1.1pt);
\draw[lold] (b3p.center) -- (c3);
\draw[lnew] (b3p.center) -- (x3.north west);
\llink{x3}{-- (c3.south)}
\llink{c3}{-- (d3)}
\lnil{d3}
\end{tikzpicture}
\diagcaption{在 $p$（$160$）後面插入 $x$（$180$）。 橘色是這一步新設的指標，灰色虛線
是被蓋掉的舊指標。 先讓 $x$ 接住後半段，再讓 $p$ 指向 $x$。}
\label{fig:ll-insert}
\end{center}

刪除也有一個順序：先把要刪的節點記在變數 $x$ 裡，接好前後，最後才釋放 $x$。 如果
先釋放，$x.\mathit{next}$ 就讀不到了。

\begin{examnote}
追蹤指標的題目，一步畫一張圖：先畫操作前的串列，每執行一行指定，就把那一個指標的
新箭頭畫上、舊箭頭劃掉。 題目問「執行完後串列是什麼」時，從 $\mathit{head}$ 沿著
箭頭重新走一次再寫答案，不要憑印象寫。 有些節點會變成沒有人指著，答案裡要寫出它們
已經不在串列上。
\end{examnote}

\subsection{依序插入與指標的指標}

報名名單要依學號排好。 插入新的學號 $k$ 時，從頭走到第一個 $\mathit{key} \ge k$
的節點 $q$，把新節點接在 $q$ 的前一個節點 $p$ 後面。 單向串列沒辦法從 $q$ 回頭找
$p$，所以走的時候讓 $p$ 跟在 $q$ 後面一格。

\begin{pseudo}[label=alg:ll-sorted-insert]{依 $\mathit{key}$ 由小到大插入}
\Procedure{Sorted-Insert}{$L, x$}
    \State $p \gets \NIL$ \Comment{$p$ 永遠跟在 $q$ 後面一格}
    \State $q \gets L.\mathit{head}$
    \While{$q \ne \NIL$ \AND $q.\mathit{key} < x.\mathit{key}$}
        \State $p \gets q$
        \State $q \gets q.\mathit{next}$
    \EndWhile
    \State \Call{List-Insert-After}{$L, p, x$}
\EndProcedure
\end{pseudo}

為什麼插完之後串列還是排好的？ 用第~1~章的迴圈不變量：每次檢查 \textbf{while} 的
條件時，從 $\mathit{head}$ 到 $p$（含 $p$）的每個節點都滿足 $\mathit{key} <
x.\mathit{key}$，而且 $q$ 是 $p$ 的下一個節點（$p = \NIL$ 時 $q$ 是第一個節點）。
一開始 $p = \NIL$，前一段是空的，成立；每一輪 $q.\mathit{key} < x.\mathit{key}$ 才
往前走，$p$ 移到 $q$ 之後仍然成立。 迴圈停下時，$q = \NIL$ 或 $q.\mathit{key} \ge
x.\mathit{key}$。 原本的串列排好了，$p$ 之前的都不大於 $p$，所以把 $x$ 放在 $p$ 和
$q$ 中間，左邊都比 $x$ 小、右邊都不比 $x$ 小，串列仍然排好。 最壞情況走完整條串
列，$\Theta(n)$。

虛擬碼裡有兩個地方要分開處理「在最前面」：$p = \NIL$ 的初值，以及
\proc{List-Insert-After} 的第一個分支。 C 有一個寫法可以把它們合成一種：不記「前一
個節點」，改記「指向 $q$ 的那個指標欄位在哪裡」。 $q$ 是第一個節點時，指著它的是
$\mathit{head}$ 這個變數；否則是前一個節點的 \texttt{next} 欄位。 兩者的型別都是
\texttt{Node *}，所以它們的位址都是 \texttt{Node **}，可以用同一個變數
\texttt{link} 記。

\begin{cbookrecap}[第~14~章 14.5 節、第~19~章 19.9、19.10 節]
\texttt{p->next} 是 \texttt{(*p).next} 的簡寫：先沿著指標 \texttt{p} 找到結構，再
取它的 \texttt{next} 欄位。

C 的參數都是傳值。 函式要改呼叫端的 \texttt{Node *head}，必須收 \texttt{Node **}，
呼叫時傳 \texttt{\&head}，在函式裡寫 \texttt{*head = ...}。 同樣地，
\texttt{\&p->next} 是節點 \texttt{p} 的 \texttt{next} 欄位的位址；拿到它，就能改
「\texttt{p} 的下一個是誰」。 第~19~章用這個寫法把「刪第一個」和「刪中間」寫成同一
段程式，這一節用同樣的寫法做依序插入。
\end{cbookrecap}

\begin{center}
\begin{tikzpicture}
\node[lname, draw=black!60, minimum width=1.0cm, minimum height=0.65cm] (hd) at (0,0) {};
\node[lname, above=1pt] at (hd.north) {head};
\lnode{a}{1.8,0}{115}
\lnode{c}{3.8,0}{207}
\lnode{d}{5.8,0}{342}
\fill[black!70] (hd.center) circle (1.1pt);
\draw[lnk] (hd.center) -- (a);
\llink{a}{-- (c)}
\llink{c}{-- (d)}
\lnil{d}
\node[lname, text=memStroke] (lk1) at (0,-1.4) {link（開始）};
\draw[lnew] (lk1) -- (hd.south);
\node[lname, text=memStroke] (lk2) at (2.4,-1.4) {link（停下）};
\draw[lnew] (lk2) -- (ap.south);
\end{tikzpicture}
\diagcaption{插入 $160$ 時 \texttt{link} 的位置。 它一開始指向 \texttt{head} 這個
變數；$115 < 160$，往前一格，改指向 $115$ 的 \texttt{next} 欄位；那個欄位指著
$207 \ge 160$，停下。 新節點就接在 \texttt{*link} 這個位置。}
\label{fig:ll-link}
\end{center}

\texttt{link} 停下時，\texttt{*link} 就是虛擬碼的 $q$。 插入只要兩行：
\texttt{p->next = *link;} 讓新節點接住 $q$，\texttt{*link = p;} 讓原本指著 $q$ 的
那個指標（不管是 \texttt{head} 還是某個 \texttt{next}）改指新節點。 刪除也一樣：
\texttt{*link = dead->next;}。 下面的程式用它管理報名名單。 為了圖好畫，學號只寫
末三碼。

\begin{example}[label=ex:ll-signup]{c}{(signup.c) 依學號排好的報名名單}
#include <stdio.h>
#include <stdlib.h>

typedef struct Node {
    int id;                 // 學號（這裡只用末三碼）
    struct Node *next;
} Node;

// 回傳「指向第一個 id >= x 的節點」的那個指標欄位的位址，
// 可能是 head 本身，也可能是某個節點的 next 欄位。
Node **find_link(Node **head, int x)
{
    Node **link = head;
    while (*link != NULL && (*link)->id < x)
        link = &(*link)->next;
    return link;
}

// 依學號由小到大插入。成功回傳 1，已經在名單上回傳 0，
// 要不到記憶體回傳 -1。
int add(Node **head, int x)
{
    Node **link = find_link(head, x);
    if (*link != NULL && (*link)->id == x)
        return 0;
    Node *p = malloc(sizeof *p);
    if (p == NULL)
        return -1;
    p->id = x;
    p->next = *link;        // 先接住後面那一段
    *link = p;              // 再讓前面指向新節點
    return 1;
}

// 刪除學號 x。成功回傳 1，找不到回傳 0。
int drop(Node **head, int x)
{
    Node **link = find_link(head, x);
    if (*link == NULL || (*link)->id != x)
        return 0;
    Node *dead = *link;
    *link = dead->next;     // 前面直接跳過 dead
    free(dead);
    return 1;
}

void print_list(const Node *p)
{
    printf("list:");
    for (; p != NULL; p = p->next)
        printf(" %d", p->id);
    printf("\n");
}

int main(void)
{
    Node *head = NULL;
    int join[] = {207, 115, 342, 160, 115, 101};
    int quit[] = {115, 300, 342};
    for (int i = 0; i < 6; i++) {
        int r = add(&head, join[i]);
        if (r < 0)
            break;
        printf("add %d: %s\n", join[i], r ? "ok" : "already in");
    }
    print_list(head);
    for (int i = 0; i < 3; i++)
        printf("drop %d: %s\n", quit[i],
               drop(&head, quit[i]) ? "ok" : "not found");
    print_list(head);
    while (head != NULL)    // 全部釋放：每次刪掉第一個
        drop(&head, head->id);
    return 0;
}
\end{example}
\iosample{}{add 207: ok \\ add 115: ok \\ add 342: ok \\ add 160: ok \\ add 115: already in \\ add 101: ok \\ list: 101 115 160 207 342 \\ drop 115: ok \\ drop 300: not found \\ drop 342: ok \\ list: 101 160 207}

這支程式有幾點可以對照前面的說明：

\begin{itemize}
\item 加入 $115$ 時它插在最前面，$342$ 插在最後面，$160$ 插在中間，三種情況用的是
同一段程式。 刪除 $115$ 時它是第一個節點，$342$ 是最後一個，也一樣。
\item \texttt{find\_link} 只用 \texttt{<} 往前走，停下時 \texttt{*link} 是第一個
\texttt{id >= x} 的節點，所以只要再比一次 \texttt{==} 就知道 \texttt{x} 在不在名
單上。 重複報名、取消不存在的學號都在這裡擋掉。
\item 程式最後一直刪第一個節點，直到名單變空，所有節點都釋放了。 \texttt{drop} 先
把 \texttt{dead->next} 接好才 \texttt{free(dead)}；順序反過來，就是讀了已經釋放的記
憶體，AddressSanitizer（《C 程式設計基礎》第~19~章 19.5 節）會直接報錯。
\end{itemize}

下表是加入 $160$ 時 \texttt{find\_link} 的過程（這時名單是 $115, 207, 342$）。

\begin{center}
\begin{tabular}{clll}
\toprule
檢查 & \texttt{link} 指向 & \texttt{*link} 指向的節點 & 結果 \\
\midrule
1 & \texttt{head} & $115$ & $115 < 160$，往前 \\
2 & $115$ 的 \texttt{next} & $207$ & $207 \ge 160$，停下 \\
\bottomrule
\end{tabular}
\end{center}

插入之後 $115$ 的 \texttt{next} 指向新節點 $160$，$160$ 的 \texttt{next} 指向
$207$。 每次操作最多走過整條名單，所以 $n$ 個人的名單每次報名或取消都是 $O(n)$。

\section{帶頭節點、環狀與雙向串列}
\label{sec:ll-variants}

上一節的程式裡，「在最前面」一直是要另外想的情況。 這一節的三種變形各自解決單向串
列的一個不方便：頭節點去掉特殊情況，環狀串列讓頭尾都能 $\Theta(1)$ 操作，雙向串列
讓節點知道前一個是誰。

\subsection{頭節點}

\textbf{頭節點}（header node，也叫 dummy node）是放在串列最前面、不存資料的節點。
串列改成記住頭節點 $h$，真正的第一個元素是 $h.\mathit{next}$；空串列時
$h.\mathit{next} = \NIL$，但 $h$ 本身一直存在。 有了它，每個資料節點前面都有一個
節點，\proc{List-Insert-After} 和 \proc{List-Delete-After} 的 $p$ 永遠不是 \NIL，
虛擬碼~\ref{alg:ll-insert} 第 2--5 行的特殊情況可以整段刪掉，只剩兩行指定。

\begin{center}
\begin{tikzpicture}
\node[lname] (h) at (-0.3,0) {h};
\lnode[lhd]{z}{0.8,0}{}
\lnode{a}{2.8,0}{12}
\lnode{b}{4.8,0}{31}
\lnil{b}
\draw[lnk] (h) -- (z);
\llink{z}{-- (a)}
\llink{a}{-- (b)}
\node[lname] (h2) at (7.4,0) {h};
\lnode[lhd]{e}{8.5,0}{}
\lnil{e}
\draw[lnk] (h2) -- (e);
\node[font=\small] at (3.0,-0.8) {兩個元素};
\node[font=\small] at (8.8,-0.8) {空串列};
\end{tikzpicture}
\diagcaption{帶頭節點的串列。 灰底的頭節點不存資料；空串列也有頭節點。}
\label{fig:ll-header}
\end{center}

代價是多一個節點的空間，而且走訪時要從 $h.\mathit{next}$ 開始。 C 程式常把頭節點
宣告成函式裡的區域變數 \texttt{Node head}，用 \texttt{\&head} 當「最後一個節點」開始
往後接，最後回傳 \texttt{head.next}；這樣接第一個節點時不用判斷串列是不是空的。
第~\ref{sec:ll-poly}~節的多項式加法和第~\ref{sec:ll-ops}~節的合併都用這個寫法。 它
和上一節的 \texttt{Node **link} 解決的是同一件事，選一個自己寫得順的就好。

\subsection{環狀串列}

把最後一個節點的 $\mathit{next}$ 接回第一個節點，就是\textbf{環狀串列}（circular
list）。 環上沒有 \NIL，任何一個節點往下走都會繞回自己。

環狀串列通常記的是\textbf{最後一個}節點 $R.\mathit{last}$。 第一
個節點就是 $R.\mathit{last}.\mathit{next}$，一步就拿得到；而最後一個節點正是「最前
面」的前一個節點，所以在最前面和在最後面插入，都是接在 $R.\mathit{last}$ 後面，差
別只在插完之後 $R.\mathit{last}$ 要不要移到新節點。

\begin{pseudo}[label=alg:ll-ring-push]{在環狀串列的最後面插入}
\Procedure{Ring-Push-Back}{$R, x$}
    \If{$R.\mathit{last} = \NIL$} \Comment{空的環：自己接自己}
        \State $x.\mathit{next} \gets x$
    \Else
        \State $x.\mathit{next} \gets R.\mathit{last}.\mathit{next}$
        \State $R.\mathit{last}.\mathit{next} \gets x$
    \EndIf
    \State $R.\mathit{last} \gets x$
\EndProcedure
\end{pseudo}

把第 7 行改成「只有原本是空的環時才做」，就是在最前面插入。 兩者都是
$\Theta(1)$。 如果只記第一個節點，在最前面插入時要改最後一個節點的
$\mathit{next}$，得先繞一圈找到它，反而變成 $\Theta(n)$。

走訪環狀串列不能用「走到 \NIL 為止」當停止條件，那會永遠走不完。 正確的寫法是從
第一個節點出發，處理完最後一個（$x = R.\mathit{last}$）就停，或是用
\textbf{do-while}：先處理再前進，前進後回到起點就停。

\begin{center}
\begin{tikzpicture}[lkey/.append style={minimum width=1.1cm}]
\lnode{a}{0,0}{Ming}
\lnode{b}{2.4,0}{Hua}
\lnode{c}{4.8,0}{Ting}
\llink{a}{-- (b)}
\llink{b}{-- (c)}
\fill[black!70] (cp.center) circle (1.1pt);
\draw[lnk, rounded corners=6pt] (cp.center) -- ++(0.5,0) -- ++(0,-0.75) -- ($(a.south)+(0,-0.42)$) -- (a.south);
\node[lname, text=memStroke] (lst) at (5.05,1.1) {last};
\draw[lnew] (lst) -- (c.north);
\end{tikzpicture}
\diagcaption{三個人的值日環。 \texttt{last} 指向 Ting，今天值日的是
\texttt{last->next}，也就是 Ming。}
\label{fig:ll-ring}
\end{center}

下面的例子是實驗室的值日輪替：每天由環上的下一位值日；新加入的人排在這一輪的最後
（也就是下一個要值日的人的前一位）；有人離開就從環上拿掉。 「過一天」只要把
\texttt{last} 往前移一格，整個環就轉了一格，不用搬任何節點。

\begin{example}[label=ex:ll-duty]{c}{(duty\_ring.c) 用環狀串列排值日}
#include <stdio.h>
#include <stdlib.h>
#include <string.h>

typedef struct Member {
    char name[16];
    struct Member *next;
} Member;

// 環只記 last：last 是這一輪的最後一位，last->next 是今天值日的人。
// 新成員接在 last 後面，再讓 last 指向他，也就是排到這一輪的最後。
int join(Member **last, const char *name)
{
    Member *p = malloc(sizeof *p);
    if (p == NULL)
        return 0;
    snprintf(p->name, sizeof p->name, "%s", name);
    if (*last == NULL) {
        p->next = p;                // 只有一個人：自己接自己
    } else {
        p->next = (*last)->next;
        (*last)->next = p;
    }
    *last = p;
    return 1;
}

// 名字是 name 的人退出。找到回傳 1，找不到回傳 0。
int leave(Member **last, const char *name)
{
    if (*last == NULL)
        return 0;
    Member *prev = *last;
    Member *cur = prev->next;       // 從今天值日的人開始繞一圈
    do {
        if (strcmp(cur->name, name) == 0) {
            if (cur == prev)            // 環上只剩他一個
                *last = NULL;
            else {
                prev->next = cur->next;
                if (cur == *last)       // 刪掉的是最後一位
                    *last = prev;
            }
            free(cur);
            return 1;
        }
        prev = cur;
        cur = cur->next;
    } while (prev != *last);
    return 0;
}

int main(void)
{
    Member *last = NULL;
    char op, name[16];
    int day = 0;
    while (scanf(" %c", &op) == 1) {
        if (op == 'N') {            // 過一天：環轉一格
            if (last == NULL)
                printf("day %d: nobody\n", ++day);
            else {
                printf("day %d: %s\n", ++day, last->next->name);
                last = last->next;
            }
        } else if (scanf("%15s", name) == 1) {
            if (op == 'J' && !join(&last, name))
                break;
            if (op == 'L' && !leave(&last, name))
                printf("%s is not on duty list\n", name);
        }
    }
    while (last != NULL)            // 全部釋放
        leave(&last, last->next->name);
    return 0;
}
\end{example}
\iosample{J Ming \\ J Hua \\ J Ting \\ N \\ N \\ J Wei \\ N \\ L Hua \\ N \\ N \\ L Ting \\ L Ming \\ L Ting \\ N \\ N \\ L Wei \\ N}{day 1: Ming \\ day 2: Hua \\ day 3: Ting \\ day 4: Ming \\ day 5: Wei \\ Ting is not on duty list \\ day 6: Wei \\ day 7: Wei \\ day 8: nobody}

\texttt{J} 是加入，\texttt{L} 是離開，\texttt{N} 是過一天。 Wei 是在第 2 天之後加入
的，那時 \texttt{last} 是 Hua，下一個輪到 Ting。 Wei 接在 Hua 後面、Ting 前面，輪替
順序變成 Ting、Ming、Hua、Wei。 Hua 在第 3 天之後就離開了，所以第 4 天是 Ming，第 5
天輪到 Wei。 環上只剩 Wei 時，他的 \texttt{next} 指向自己，每天都是他。

\texttt{leave} 有三種情況要分開想：刪掉之後環變空了、刪掉的是 \texttt{last}（
\texttt{last} 要退回前一個）、其他。 它的迴圈用 \textbf{do-while}，停止條件是「剛檢
查完的是 \texttt{last}」，繞完一圈就停。 程式結束前，每次讓今天值日的人
（\texttt{last->next}）離開，\texttt{leave} 第一個就找到他，一次 $\Theta(1)$，直到環
變空。 第~\ref{sec:ll-avail}~節會看到，如果節點是還給自己管理的節點池，整個環可以一
次 $\Theta(1)$ 還回去。

\subsection{雙向串列}

單向串列拿到一個節點 $x$，沒有辦法 $\Theta(1)$ 把它刪掉，因為改的是前一個節點的
$\mathit{next}$，而前一個節點只能從頭找。 \textbf{雙向串列}（doubly linked list）
的每個節點多一個 $\mathit{prev}$ 指向前一個節點，兩個方向都能走，刪除任何已知的節
點都是 $\Theta(1)$。

雙向串列最好寫的形式是\textbf{帶頭節點的環狀雙向串列}：頭節點 $h$ 的
$\mathit{next}$ 是第一個元素，$\mathit{prev}$ 是最後一個元素；空串列時
$h.\mathit{next} = h.\mathit{prev} = h$。 這樣每個資料節點的前後都一定有節點（可能
是 $h$），插入、刪除完全沒有特殊情況。

\begin{center}
\begin{tikzpicture}[
    dn/.style={draw=black!60, minimum width=0.85cm, minimum height=0.65cm, font=\ttfamily\small},
]
\node[dn, fill=black!8] (h) at (0,0) {h};
\node[dn] (a) at (2.2,0) {12};
\node[dn] (b) at (4.4,0) {31};
\node[dn] (c) at (6.6,0) {47};
\foreach \u/\v in {h/a, a/b, b/c} {
    \draw[lnk] ([yshift=4pt]\u.east) -- ([yshift=4pt]\v.west);
    \draw[lnk] ([yshift=-4pt]\v.west) -- ([yshift=-4pt]\u.east);
}
\draw[lnk, rounded corners=6pt] (c.north) -- ++(0,0.55) -| (h.north);
\draw[lnk, rounded corners=6pt] ([xshift=-6pt]h.south) -- ++(0,-0.55) -| ([xshift=-6pt]c.south);
\node[font=\scriptsize, text=black!60] at (3.3,1.12) {最後一個的 next 是 h};
\node[font=\scriptsize, text=black!60] at (3.3,-1.12) {h 的 prev 是最後一個};
\end{tikzpicture}
\diagcaption{帶頭節點的環狀雙向串列。 上排箭頭是 $\mathit{next}$，下排是
$\mathit{prev}$。}
\label{fig:ll-doubly}
\end{center}

\begin{pseudo}[label=alg:ll-dinsert]{雙向串列：把 $x$ 接在 $p$ 後面}
\Procedure{DList-Insert-After}{$p, x$}
    \State $x.\mathit{prev} \gets p$
    \State $x.\mathit{next} \gets p.\mathit{next}$
    \State $p.\mathit{next}.\mathit{prev} \gets x$
    \State $p.\mathit{next} \gets x$
\EndProcedure
\end{pseudo}

\begin{pseudo}[label=alg:ll-dremove]{雙向串列：拿掉節點 $x$}
\Procedure{DList-Remove}{$x$}
    \State $x.\mathit{prev}.\mathit{next} \gets x.\mathit{next}$
    \State $x.\mathit{next}.\mathit{prev} \gets x.\mathit{prev}$
\EndProcedure
\end{pseudo}

插入要改四個指標：新節點自己的兩個、前一個的 $\mathit{next}$、後一個的
$\mathit{prev}$。 順序的限制只有一條：第 5 行會蓋掉 $p.\mathit{next}$，而第 3、4 行
都要用到舊的 $p.\mathit{next}$，所以這兩行要在第 5 行之前做。 第 4 行若改寫成
$x.\mathit{next}.\mathit{prev} \gets x$，限制就換成「要在第 3 行之後」。 在頭節點
$h$ 後面插入是放到最前面，在 $h.\mathit{prev}$ 後面插入是放到
最後面。 刪除只改兩個指標，被刪的 $x$ 自己的指標不用動。

\begin{center}
\begin{tikzpicture}[
    dn/.style={draw=black!60, minimum width=0.85cm, minimum height=0.65cm, font=\ttfamily\small},
]
\node[dn] (p) at (0,0) {31};
\node[dn] (q) at (4.4,0) {47};
\node[dn, draw=memStroke, thick, fill=memFill] (x) at (2.2,-2.0) {40};
\draw[lold] ([yshift=4pt]p.east) -- ([yshift=4pt]q.west);
\draw[lold] ([yshift=-4pt]q.west) -- ([yshift=-4pt]p.east);
\draw[lnew] (x.west) to[bend left=25] node[stepno, pos=0.5, left=2pt] {2} (p.south west);
\draw[lnew] (x.north east) -- node[stepno, pos=0.5, left=2pt] {3} (q.south west);
\draw[lnew] (q.south east) to[bend left=25] node[stepno, pos=0.5, right=2pt] {4} (x.east);
\draw[lnew] (p.south east) -- node[stepno, pos=0.5, right=2pt] {5} (x.north west);
\node[font=\scriptsize, text=memStroke] at (-0.6,0.5) {$p$};
\end{tikzpicture}
\diagcaption{在 $p$（$31$）後面插入 $x$（$40$）。 圓圈裡的數字是虛擬碼~\ref{alg:ll-dinsert}
的行號，灰色虛線是被取代的兩個舊指標。}
\label{fig:ll-dinsert}
\end{center}

翻成 C 時，節點多一個 \texttt{struct DNode *prev} 欄位，四行指定一行對一行照抄（例如
\texttt{x->prev = p;}）；空串列由 \texttt{h->next = h->prev = h} 開始。
\proc{DList-Remove} 不負責釋放，拿下來的節點可以 \texttt{free}，也可以接到別的地方，
例如再接到 $h$ 後面，就是「移到最前面」。

\subsection{幾種串列的比較}

下表整理本節三種結構和單向串列的成本。 「單向」只記 $\mathit{head}$，「單向＋tail」
另外記最後一個節點，「單向環狀」只記 $\mathit{last}$，「雙向環狀」帶頭節點。 「已知
節點 $x$」是說手上已經有指向 $x$ 的指標，不包含找到它的時間。

\begin{center}
{\small
\begin{tabular}{lcccc}
\toprule
操作 & 單向 & 單向＋tail & 單向環狀 & 雙向環狀 \\
\midrule
在最前面插入 & $\Theta(1)$ & $\Theta(1)$ & $\Theta(1)$ & $\Theta(1)$ \\
在最後面插入 & $\Theta(n)$ & $\Theta(1)$ & $\Theta(1)$ & $\Theta(1)$ \\
刪除第一個 & $\Theta(1)$ & $\Theta(1)$ & $\Theta(1)$ & $\Theta(1)$ \\
刪除最後一個 & $\Theta(n)$ & $\Theta(n)$ & $\Theta(n)$ & $\Theta(1)$ \\
刪除已知節點 $x$ & $\Theta(n)$ & $\Theta(n)$ & $\Theta(n)$ & $\Theta(1)$ \\
每個節點的指標數 & 1 & 1 & 1 & 2 \\
\bottomrule
\end{tabular}}
\end{center}

「刪除最後一個」在單向串列裡都是 $\Theta(n)$，即使記了 \texttt{tail} 也一樣：刪掉
之後 \texttt{tail} 要改成倒數第二個，而倒數第二個只能從頭找。 第~6~章用串列做佇列
時，會選「在最後面插入、刪除第一個」這兩個在單向串列上都便宜的操作。

\begin{examnote}
寫串列的程式題時，寫完後拿四種情況各檢查一次：空串列、只有一個節點、操作的是第一
個節點、操作的是最後一個節點。 大部分錯誤（讀到 \NIL 的欄位、\texttt{head} 或
\texttt{last} 沒有更新）都出在這四種情況。 有頭節點或用 \texttt{Node **} 時，可以
在答案裡說明為什麼這些情況不必另外處理。
\end{examnote}

\section{用串列表示多項式與稀疏矩陣}
\label{sec:ll-poly}

第~4~章用陣列存多項式和稀疏矩陣：多項式可以存全部的係數，也可以只存非零項；稀疏
矩陣用三元組 $(\text{列}, \text{行}, \text{值})$ 的陣列。 只存非零項省空間，但項數
會變的時候（相加後有些項相消、相乘後項數變多），陣列要預先開多大很難說，插入一項還
得挪動後面的項。 改用串列，每一項是一個節點，要多少就配置多少。

\subsection{多項式}

多項式的每一個非零項存成一個節點，欄位是係數 $\mathit{coef}$、次方 $\mathit{exp}$
和 $\mathit{next}$，節點依次方由大到小串起來。 例如
\[
A(x) = 5x^{12} - 3x^{7} + x^{2} + 8, \qquad B(x) = 2x^{9} + 3x^{7} - x^{2} + 4x
\]
存成下圖的兩條串列，每個節點寫成「係數, 次方」。

\begin{center}
\begin{tikzpicture}[lkey/.append style={minimum width=1.05cm}]
\node[lname] (ha) at (-0.4,0) {A};
\lnode{a1}{0.8,0}{5,12}
\lnode{a2}{3.0,0}{-3,7}
\lnode{a3}{5.2,0}{1,2}
\lnode{a4}{7.4,0}{8,0}
\draw[lnk] (ha) -- (a1);
\llink{a1}{-- (a2)}
\llink{a2}{-- (a3)}
\llink{a3}{-- (a4)}
\lnil{a4}
\node[lname] (hb) at (-0.4,-1.1) {B};
\lnode{b1}{0.8,-1.1}{2,9}
\lnode{b2}{3.0,-1.1}{3,7}
\lnode{b3}{5.2,-1.1}{-1,2}
\lnode{b4}{7.4,-1.1}{4,1}
\draw[lnk] (hb) -- (b1);
\llink{b1}{-- (b2)}
\llink{b2}{-- (b3)}
\llink{b3}{-- (b4)}
\lnil{b4}
\end{tikzpicture}
\diagcaption{多項式 $A$、$B$ 的串列，依次方由大到小。 係數是 $0$ 的項不存。}
\label{fig:ll-poly}
\end{center}

相加的做法和合併兩條排好的串列一樣：$a$、$b$ 各指著一條串列目前的項，次方大的那一
項先抄到結果；次方一樣就把係數加起來，和不是 $0$ 才抄，兩邊一起前進。 有一邊走完
了，另一邊剩下的項全部照抄。

\begin{pseudo}[label=alg:ll-poly-add]{多項式相加}
\Function{Poly-Add}{$a, b$} \Comment{$a$、$b$ 指向兩個多項式的第一項}
    \State $h \gets$ 新的頭節點；$t \gets h$ \Comment{$t$ 是結果的最後一項}
    \While{$a \ne \NIL$ \AND $b \ne \NIL$}
        \If{$a.\mathit{exp} > b.\mathit{exp}$}
            \State $t \gets \Call{Attach}{t, a.\mathit{coef}, a.\mathit{exp}}$
            \State $a \gets a.\mathit{next}$
        \ElsIf{$a.\mathit{exp} < b.\mathit{exp}$}
            \State $t \gets \Call{Attach}{t, b.\mathit{coef}, b.\mathit{exp}}$
            \State $b \gets b.\mathit{next}$
        \Else
            \State $s \gets a.\mathit{coef} + b.\mathit{coef}$
            \If{$s \ne 0$}
                \State $t \gets \Call{Attach}{t, s, a.\mathit{exp}}$
            \EndIf
            \State $a \gets a.\mathit{next}$；$b \gets b.\mathit{next}$
        \EndIf
    \EndWhile
    \State 把 $a$ 或 $b$ 剩下的項依序 \proc{Attach} 到 $t$ 後面
    \State \Return $h.\mathit{next}$
\EndFunction
\end{pseudo}

\proc{Attach}$(t, c, e)$ 配置一個新節點、接在 $t$ 後面，並回傳這個新節點（新的最
後一項）。 結果用頭節點 $h$ 起頭，第一次 \proc{Attach} 就不用另外處理。

\textbf{成本}。 設 $A$ 有 $m$ 項、$B$ 有 $n$ 項。 迴圈每一輪比較一次次方，至少讓
$a$、$b$ 其中一個前進一格，所以迴圈最多 $m + n$ 輪；更仔細地說，最後一個被處理的
項不會在迴圈裡被比較，比較次數最多 $m + n - 1$。 剩下的項照抄，總時間
$\Theta(m + n)$，和多項式的次數無關。 如果用長度 $d + 1$ 的係數陣列（$d$ 是次方），
相加要 $\Theta(d)$；$A(x) = x^{1000} + 1$ 這種項數少、次方高的多項式，串列快得多。

\begin{example}[label=ex:ll-poly]{c}{(poly\_add.c) 多項式相加}
#include <stdio.h>
#include <stdlib.h>

typedef struct Term {
    int coef;               // 係數，不會是 0
    int exp;                // 次方，由大到小排
    struct Term *next;
} Term;

// 在 tail 後面接上一個新項，回傳新的尾巴；要不到記憶體回傳 NULL。
Term *attach(Term *tail, int coef, int exp)
{
    Term *t = malloc(sizeof *t);
    if (t == NULL)
        return NULL;
    t->coef = coef;
    t->exp = exp;
    t->next = NULL;
    tail->next = t;
    return t;
}

// 從陣列 c[0..m-1]、e[0..m-1] 建一個多項式。
Term *build(const int *c, const int *e, int m)
{
    Term head = {0, 0, NULL};       // 頭節點：只用它的 next
    Term *t = &head;
    for (int i = 0; i < m && t != NULL; i++)
        t = attach(t, c[i], e[i]);
    return head.next;
}

// 回傳 a + b，a、b 不變。*cmp 記錄比較了幾次次方。
Term *poly_add(const Term *a, const Term *b, int *cmp)
{
    Term head = {0, 0, NULL};
    Term *t = &head;
    *cmp = 0;
    while (a != NULL && b != NULL && t != NULL) {
        (*cmp)++;
        if (a->exp > b->exp) {
            t = attach(t, a->coef, a->exp);
            a = a->next;
        } else if (a->exp < b->exp) {
            t = attach(t, b->coef, b->exp);
            b = b->next;
        } else {
            if (a->coef + b->coef != 0)     // 係數相消就不留這一項
                t = attach(t, a->coef + b->coef, a->exp);
            a = a->next;
            b = b->next;
        }
    }
    for (; a != NULL && t != NULL; a = a->next)     // 剩下的照抄
        t = attach(t, a->coef, a->exp);
    for (; b != NULL && t != NULL; b = b->next)
        t = attach(t, b->coef, b->exp);
    return head.next;
}

// 每一項印成 (係數,次方)，然後釋放整條串列。
void print_and_free(const char *name, Term *p)
{
    printf("%s:", name);
    while (p != NULL) {
        Term *next = p->next;
        printf(" (%d,%d)", p->coef, p->exp);
        free(p);
        p = next;
    }
    printf("\n");
}

int main(void)
{
    int ca[] = {5, -3, 1, 8}, ea[] = {12, 7, 2, 0};
    int cb[] = {2, 3, -1, 4}, eb[] = {9, 7, 2, 1};
    Term *a = build(ca, ea, 4);
    Term *b = build(cb, eb, 4);
    int cmp;
    Term *c = poly_add(a, b, &cmp);
    print_and_free("A", a);
    print_and_free("B", b);
    print_and_free("A+B", c);
    printf("exponent comparisons = %d\n", cmp);
    return 0;
}
\end{example}
\iosample{}{A: (5,12) (-3,7) (1,2) (8,0) \\ B: (2,9) (3,7) (-1,2) (4,1) \\ A+B: (5,12) (2,9) (4,1) (8,0) \\ exponent comparisons = 5}

$A$、$B$ 的係數和次方寫在 \texttt{main} 的陣列裡，次方由大到小。 輸出把每一項印成
（係數, 次方），結果是 $A + B = 5x^{12} + 2x^{9} + 4x + 8$。

逐輪追蹤（「抄」是把那一項接到結果的最後面）：

\begin{center}
\begin{tabular}{cccl}
\toprule
比較 & $a$ 目前的項 & $b$ 目前的項 & 動作 \\
\midrule
1 & $5x^{12}$ & $2x^{9}$ & $12 > 9$，抄 $5x^{12}$ \\
2 & $-3x^{7}$ & $2x^{9}$ & $7 < 9$，抄 $2x^{9}$ \\
3 & $-3x^{7}$ & $3x^{7}$ & 係數和 $0$，不抄，兩邊都前進 \\
4 & $x^{2}$ & $-x^{2}$ & 係數和 $0$，不抄，兩邊都前進 \\
5 & $8$ & $4x$ & $0 < 1$，抄 $4x$；$b$ 走完了 \\
-- & $8$ & -- & 迴圈結束，抄 $8$ \\
\bottomrule
\end{tabular}
\end{center}

$m = n = 4$，比較了 $5$ 次，少於上限 $m + n - 1 = 7$，因為有兩次是一次消掉兩項。

多項式乘法可以拿 $A$ 的每一項去乘 $B$ 的每一項，得到 $mn$ 個乘積項，再一項一項加
進結果串列。 結果最多 $mn$ 項，每加一項要在結果裡找位置，最壞情況要
$O((mn)^2)$；改成每次把「$A$ 的一項乘上整個 $B$」得到的一條排好的串列，用
\proc{Poly-Add} 加進結果。 第 $i$ 次相加時結果最多 $(i-1)n$ 項，這次相加是 $O(in)$，
$m$ 次合計 $O(m^2 n)$。

\subsection{稀疏矩陣：正交串列}

稀疏矩陣如果常常要改（把某一格改成 $0$、或把 $0$ 改成非零），三元組陣列每次都要挪
動，這時可以改用\textbf{正交串列}（orthogonal list，也叫十字串列）。 這一小節只講
結構和成本，不寫程式。

每個非零元素是一個節點，有五個欄位：列號 $\mathit{row}$、行號 $\mathit{col}$、值
$\mathit{val}$，以及兩個指標 $\mathit{right}$（同一列的下一個非零元素）和
$\mathit{down}$（同一行的下一個非零元素）。 每一列有一個列頭節點，把這一列的非零元
素依行號串成一條環狀串列；每一行也有一個行頭節點，把這一行的非零元素依列號串成環。
一個節點同時在一條「橫的」串列和一條「直的」串列上，所以叫做正交。 下圖左邊是矩陣
$M$，右邊是它的正交串列。

\begin{center}
$M = \begin{pmatrix}
0 & 0 & 3 & 0 \\
8 & 0 & 0 & 0 \\
0 & 0 & 0 & 0 \\
0 & 6 & 0 & 2
\end{pmatrix}$\hspace{1.2cm}%
\begin{tikzpicture}[baseline=(current bounding box.center),
    hd/.style={draw=black!60, fill=black!8, minimum width=0.8cm, minimum height=0.6cm, font=\small},
    en/.style={draw=coreStroke, fill=coreFill, minimum width=1.25cm, minimum height=0.6cm, font=\scriptsize},
    lnk/.style={-{Stealth[length=1.6mm]}, thick, draw=black!70},
]
\foreach \j in {1,...,4} { \node[hd] (C\j) at (\j*1.9,0) {C\j}; }
\foreach \i in {1,...,4} { \node[hd] (R\i) at (0,-\i*0.95) {R\i}; }
\node[en] (e13) at (3*1.9,-1*0.95) {1,3 \quad 3};
\node[en] (e21) at (1*1.9,-2*0.95) {2,1 \quad 8};
\node[en] (e42) at (2*1.9,-4*0.95) {4,2 \quad 6};
\node[en] (e44) at (4*1.9,-4*0.95) {4,4 \quad 2};
% right 指標（同一列）
\draw[lnk] (R1) -- (e13);
\draw[lnk] (R2) -- (e21);
\draw[lnk] (R4) -- (e42);
\draw[lnk] (e42) -- (e44);
% down 指標（同一行）
\draw[lnk, draw=memStroke] (C1) -- (e21);
\draw[lnk, draw=memStroke] (C2) -- (e42);
\draw[lnk, draw=memStroke] (C3) -- (e13);
\draw[lnk, draw=memStroke] (C4) -- (e44);
\end{tikzpicture}
\diagcaption{$M$ 的正交串列。 黑色箭頭是 $\mathit{right}$，橘色箭頭是
$\mathit{down}$；每個節點寫「列, 行　值」。 每一列、每一行的最後一個節點指回自己的
頭節點，圖上沒有畫出這些指標。 第 3 列沒有非零元素，R3 的 $\mathit{right}$ 指向自
己。}
\label{fig:ll-ortho}
\end{center}

幾個成本：

\begin{itemize}
\item \textbf{空間}：$r \times c$ 的矩陣有 $t$ 個非零元素時，共 $t$ 個元素節點和
$r + c$ 個頭節點，$\Theta(t + r + c)$。 三元組陣列是 $\Theta(t)$，正交串列多了頭節
點，每個節點也比三元組多兩個指標。
\item \textbf{取 $M[i][j]$}：沿第 $i$ 列走，$O(\text{第 } i \text{ 列的非零個數})$。
\item \textbf{設定 $M[i][j]$}（新增或刪除一個節點）：在第 $i$ 列找到插入位置的前一個節
點，在第 $j$ 行也找一次，各改一個指標。 時間是兩條串列的長度，不用挪動其他元素。
\item \textbf{轉置}：只要把每個節點的 $\mathit{row}$ 和 $\mathit{col}$、
$\mathit{right}$ 和 $\mathit{down}$ 對調，列頭和行頭對調，$\Theta(t + r + c)$。
\end{itemize}

矩陣只讀不改時，三元組陣列比較省空間，也比較好寫；元素常常變動時，正交串列省下挪動
的成本。 第~21~章圖的鄰接串列，可以看成只保留「橫的」那一半的正交串列。

\section{可用空間串列}
\label{sec:ll-avail}

到目前為止，每個節點都用 \texttt{malloc} 要、用 \texttt{free} 還。 有些情況不想或
不能這樣做：沒有作業系統的小型裝置可能沒有 \texttt{malloc}；節點很小、配置很頻繁時，
每次都呼叫 \texttt{malloc} 的額外成本比節點本身還大；有些考題也會直接給兩個陣列
$\mathit{key}[\,]$、$\mathit{next}[\,]$，要你用陣列模擬串列。 這時可以一次準備好一池
節點，自己管理哪些在用、哪些空著。

空著的節點也串成一條串列，叫\textbf{可用空間串列}（available-space list，或 free
list），開頭記在 $\mathit{avail}$。 要節點就從這條串列的最前面拿一個，還節點就放回
最前面，兩者都是 $\Theta(1)$。 可用空間串列其實就是一個堆疊（第~6~章）。

\begin{pseudo}[label=alg:ll-getnode, unbreakable]{從可用空間串列拿一個節點}
\Function{Get-Node}{}
    \If{$\mathit{avail} = \NIL$}
        \State \Error{沒有可用的節點}
    \EndIf
    \State $x \gets \mathit{avail}$
    \State $\mathit{avail} \gets x.\mathit{next}$
    \State \Return $x$
\EndFunction
\end{pseudo}

\begin{pseudo}[label=alg:ll-retnode]{把節點 $x$ 還回去}
\Procedure{Ret-Node}{$x$}
    \State $x.\mathit{next} \gets \mathit{avail}$
    \State $\mathit{avail} \gets x$
\EndProcedure
\end{pseudo}

一開始所有節點都空著，把它們依序串起來就是最初的可用空間串列。 用陣列實作時，
「指標」改成陣列的編號，\NIL 用一個不會是編號的值（例如 $-1$）表示；
$x.\mathit{next}$ 寫成 $\mathit{next}[x]$。 這種寫法也叫\textbf{游標}（cursor）實
作。 一般的串列操作照寫，只是 \texttt{malloc} 換成 \proc{Get-Node}、\texttt{free}
換成 \proc{Ret-Node}；刪除時一樣要找到前一個節點，插入一樣先接後面再改前面。

下面用一個只有 $5$ 格的節點池追蹤一次。 一開始 $\mathit{avail} = 0$，
$\mathit{next}[\,] = \langle 1, 2, 3, 4, -1 \rangle$，串列是空的（$\mathit{head} =
-1$）。 「插入」都插在最前面；「$40@0$」表示 $40$ 放在第 $0$ 格。

\begin{center}
\begin{tabular}{llll}
\toprule
操作 & 用到的格子 & 串列（從 $\mathit{head}$ 走） & 可用空間串列 \\
\midrule
插入 40 & 拿第 0 格 & $40@0$ & 1, 2, 3, 4 \\
插入 25 & 拿第 1 格 & $25@1, 40@0$ & 2, 3, 4 \\
插入 63 & 拿第 2 格 & $63@2, 25@1, 40@0$ & 3, 4 \\
刪除 25 & 還第 1 格 & $63@2, 40@0$ & 1, 3, 4 \\
插入 18 & 拿第 1 格 & $18@1, 63@2, 40@0$ & 3, 4 \\
插入 77 & 拿第 3 格 & $77@3, 18@1, 63@2, 40@0$ & 4 \\
插入 52 & 拿第 4 格 & $52@4, 77@3, 18@1, 63@2, 40@0$ & （空） \\
插入 11 & 失敗 & 不變 & （空） \\
刪除 18 & 還第 1 格 & $52@4, 77@3, 63@2, 40@0$ & 1 \\
\bottomrule
\end{tabular}
\end{center}

\begin{center}
最後的陣列是\quad
\begin{tabular}{cccccc}
\toprule
編號 & 0 & 1 & 2 & 3 & 4 \\
\midrule
$\mathit{key}$ & 40 & 18 & 63 & 77 & 52 \\
$\mathit{next}$ & $-1$ & $-1$ & 0 & 2 & 3 \\
\bottomrule
\end{tabular}
\qquad $\mathit{head} = 4$，$\mathit{avail} = 1$。
\end{center}

幾件要看清楚的事：

\begin{itemize}
\item 刪除 $25$ 把第 $1$ 格還回去，它被放在可用空間串列的\textbf{最前面}，所以下一
次插入 $18$ 拿到的又是第 $1$ 格。 剛還的節點最先被重用，和堆疊一樣後進先出。
\item 五格都用完之後，插入 $11$ 失敗。 這是自己管理節點池的代價：容量要事先決定。
\item 第 $1$ 格的 $\mathit{key}$ 還是 $18$。 還回去的節點不會被清掉，只是不在串列上
了；判斷一格是不是在用，要看它在哪一條串列上，不能看裡面的值。
\end{itemize}

\begin{examnote}
用陣列模擬串列的追蹤題，每一步都寫出 $\mathit{head}$、$\mathit{avail}$ 和有變動的
$\mathit{next}[\,]$。 最後答案除了陣列內容，最好再從 $\mathit{head}$ 和
$\mathit{avail}$ 各走一次，寫出兩條串列，例如「串列：$4, 3, 2, 0$；可用：$1$」。
這兩條串列加起來要剛好是每一格各一次，可以拿來驗算。
\end{examnote}

\subsection{一次還一整條環狀串列}

要把一條有 $k$ 個節點的單向串列整條還回去，得先走到它的最後一個節點，才能把它接到
$\mathit{avail}$ 前面，$\Theta(k)$。 環狀串列記的正是最後一個節點，所以可以
$\Theta(1)$ 做完：

\enlargethispage{\baselineskip}% 讓整個虛擬碼框留在同一頁
\begin{pseudo}[label=alg:ll-ring-erase]{把整個環還給可用空間串列}
\Procedure{Ring-Erase}{$R$}
    \If{$R.\mathit{last} \ne \NIL$}
        \State $\mathit{first} \gets R.\mathit{last}.\mathit{next}$
        \State $R.\mathit{last}.\mathit{next} \gets \mathit{avail}$
        \State $\mathit{avail} \gets \mathit{first}$
        \State $R.\mathit{last} \gets \NIL$
    \EndIf
\EndProcedure
\end{pseudo}

第 4 行把環在 $\mathit{last}$ 後面剪開、接上原本的可用空間串列，第 5 行讓
$\mathit{avail}$ 從環的第一個節點開始。 做完之後，原本環上的 $k$ 個節點排在可用空間
串列的最前面，順序和在環上一樣。 和第~3~章的平攤分析放在一起看：建立一個 $k$ 個
節點的環要 $k$ 次 \proc{Get-Node}，拆掉只要 $\Theta(1)$，一連串操作的總時間只和配
置過的節點數有關。

\section{常見操作：反轉、合併、找中間}
\label{sec:ll-ops}

這一節的三個操作都只改 $\mathit{next}$，不配置新節點，額外空間 $\Theta(1)$。 它們
常常一起出現：第~14~章串列版的合併排序，就是「找中間切成兩半、各自排好、再合併」。
這一節用同一個情境示範：兩位助教各自登記了一部分同學的小考成績，名單都依學號排好；
要把它們合成一份，再切成前後兩半分到兩間教室，後半那間要從學號大的開始點名。

\subsection{反轉}

\textbf{反轉}把串列 $x_1 \to x_2 \to \cdots \to x_n$ 改成 $x_n \to \cdots \to x_1$。
做法是從頭一個一個拆下來，放到另一條串列 $d$ 的最前面。 拆之前要先記住下一個節點，
因為改掉 $h.\mathit{next}$ 之後就找不到它了。

\begin{pseudo}[label=alg:ll-reverse, unbreakable]{就地反轉}
\Function{List-Reverse}{$h$} \Comment{$h$ 指向第一個節點}
    \State $d \gets \NIL$ \Comment{已經反轉好的一段}
    \While{$h \ne \NIL$}
        \State $r \gets h.\mathit{next}$ \Comment{先記住還沒處理的部分}
        \State $h.\mathit{next} \gets d$
        \State $d \gets h$
        \State $h \gets r$
    \EndWhile
    \State \Return $d$
\EndFunction
\end{pseudo}

用 $112 \to 115 \to 117 \to 120$ 追蹤，每一列是 \textbf{while} 做完一輪之後的狀態：

\begin{center}
\begin{tabular}{clll}
\toprule
輪 & 這一輪搬的節點 & $d$ 那一段 & $h$ 那一段 \\
\midrule
開始 & -- & （空） & $112 \to 115 \to 117 \to 120$ \\
1 & $112$ & $112$ & $115 \to 117 \to 120$ \\
2 & $115$ & $115 \to 112$ & $117 \to 120$ \\
3 & $117$ & $117 \to 115 \to 112$ & $120$ \\
4 & $120$ & $120 \to 117 \to 115 \to 112$ & （空） \\
\bottomrule
\end{tabular}
\end{center}

\begin{theorem}
\label{thm:ll-reverse}
\proc{List-Reverse} 回傳反轉後的串列，時間 $\Theta(n)$，額外空間 $\Theta(1)$。
\end{theorem}

\begin{proof}
設原本的串列是 $x_1 \to \cdots \to x_n$。 不變量：第 $i$ 輪開始時（$1 \le i \le
n + 1$），$d$ 是 $x_{i-1} \to x_{i-2} \to \cdots \to x_1 \to \NIL$，$h$ 是 $x_i \to
x_{i+1} \to \cdots \to x_n \to \NIL$，而且 $x_i$ 以後的 $\mathit{next}$ 都沒有被改過
（$i = n + 1$ 時 $h = \NIL$）。

初始：$i = 1$，$d = \NIL$ 是空的一段，$h$ 是整條原本的串列，成立。

維持：第 $i$ 輪 $h = x_i \ne \NIL$。 第 4 行記下 $r = x_{i+1}$（$i = n$ 時是
\NIL）；第 5 行讓 $x_i$ 指向 $d$，$d$ 這一段變成 $x_i \to x_{i-1} \to \cdots \to
x_1$；第 6、7 行讓 $d = x_i$、$h = x_{i+1}$。 第 5 行只改了 $x_i$ 的
$\mathit{next}$，$x_{i+1}$ 以後沒有動，所以 $h$ 仍然是 $x_{i+1} \to \cdots \to x_n$。
下一輪開始時不變量對 $i + 1$ 成立。

結束：迴圈在 $h = \NIL$ 時停，也就是 $i = n + 1$。 這時 $d$ 是 $x_n \to \cdots \to
x_1 \to \NIL$，正是反轉後的串列。

迴圈跑 $n$ 輪，每輪 $\Theta(1)$；只用了 $d$、$r$ 兩個額外的指標。
\end{proof}

反轉也可以寫成遞迴：先反轉 $x_2$ 以後的部分，再把 $x_1$ 接到它的最後面。 遞迴的深
度是 $n$，要 $\Theta(n)$ 的堆疊空間，串列很長時可能用光呼叫堆疊（《C 程式設計基礎》
第~19~章 19.7 節）。 觀念題第 7 題會證明遞迴版的正確性。

\subsection{合併兩條排好的串列}

兩條依 $\mathit{key}$ 由小到大排好的串列，要合成一條排好的串列。 和多項式相加一
樣，兩邊各一個指標，每次把比較小的那個節點接到結果的最後面；不同的是這裡直接搬原本
的節點，不配置新的。

\begin{pseudo}[label=alg:ll-merge]{合併兩條排好的串列}
\Function{Merge-Lists}{$a, b$}
    \State $h \gets$ 頭節點；$t \gets h$
    \While{$a \ne \NIL$ \AND $b \ne \NIL$}
        \If{$a.\mathit{key} \le b.\mathit{key}$}
            \State $t.\mathit{next} \gets a$；$a \gets a.\mathit{next}$
        \Else
            \State $t.\mathit{next} \gets b$；$b \gets b.\mathit{next}$
        \EndIf
        \State $t \gets t.\mathit{next}$
    \EndWhile
    \If{$a \ne \NIL$}
        \State $t.\mathit{next} \gets a$ \Comment{剩下的一整段直接接上}
    \Else
        \State $t.\mathit{next} \gets b$
    \EndIf
    \State \Return $h.\mathit{next}$
\EndFunction
\end{pseudo}

兩點說明。 第一，第 4 行用 $\le$：$\mathit{key}$ 相同時先拿 $a$ 的，所以同樣
$\mathit{key}$ 的節點保持「$a$ 的在前、$b$ 的在後」。 這叫做\textbf{穩定}
（stable），第~14~章的合併排序靠這一點才是穩定的排序。 第二，某一邊用完時，另一邊
剩下的部分本來就排好了，而且都不比已經接上的節點小，整段接上只要 $\Theta(1)$，不用
一個一個搬。 這是串列比陣列方便的地方。

設兩條串列各有 $m$、$n$ 個節點。 每比較一次，就有一個節點接到結果上，而迴圈在其中
一邊用完時就停，所以比較次數介於 $\min(m, n)$ 和 $m + n - 1$ 之間（觀念題第 6 題會證
明兩個界都達得到）。 剩下的一段只花 $\Theta(1)$ 接上，所以總時間也介於
$\Theta(\min(m, n))$ 和 $\Theta(m + n)$ 之間，一般寫成 $O(m + n)$。 陣列版的合併
必須把剩下的元素一個一個抄過去，永遠是 $\Theta(m + n)$。

\subsection{找中間：快慢指標}

要把串列切成前後兩半，得先知道中間在哪裡。 直接的做法是先走一次數出長度 $n$，再走
$\lceil n/2 \rceil$ 步。 另一種做法只走一次：兩個指標一起從頭出發，\textbf{慢指
標} $s$ 每次走一步，\textbf{快指標} $f$ 每次走兩步。 $f$ 走到尾巴時，$s$ 剛好在中
間。

\begin{pseudo}[label=alg:ll-split]{切成前後兩半，前半有 $\lceil n/2 \rceil$ 個節點}
\Function{Split-Half}{$h$}
    \If{$h = \NIL$}
        \State \Return \NIL
    \EndIf
    \State $s \gets h$；$f \gets h.\mathit{next}$
    \While{$f \ne \NIL$ \AND $f.\mathit{next} \ne \NIL$}
        \State $s \gets s.\mathit{next}$
        \State $f \gets f.\mathit{next}.\mathit{next}$
    \EndWhile
    \State $b \gets s.\mathit{next}$；$s.\mathit{next} \gets \NIL$
    \State \Return $b$ \Comment{後半的開頭；前半仍從 $h$ 開始}
\EndFunction
\end{pseudo}

\begin{theorem}
\label{thm:ll-split}
設 $h$ 是 $x_1 \to \cdots \to x_n$，$n \ge 1$。 \proc{Split-Half} 的迴圈跑
$\lfloor (n-1)/2 \rfloor$ 輪，結束時 $s = x_{\lceil n/2 \rceil}$，所以前半是
$x_1, \ldots, x_{\lceil n/2 \rceil}$。 $n \ge 2$ 時 $1 \le \lceil n/2 \rceil < n$，
前後兩半都不是空的。
\end{theorem}

\begin{proof}
做完 $k$ 輪之後，$s = x_{1+k}$，$f = x_{2+2k}$（$2 + 2k > n$ 時 $f = \NIL$）。
這對 $k = 0$ 成立，每一輪 $s$ 前進一格、$f$ 前進兩格，所以一直成立。 迴圈繼續的條
件是 $x_{2+2k}$ 和 $x_{3+2k}$ 都存在，也就是 $2k + 3 \le n$。 所以迴圈跑的輪數
$K$ 是滿足 $2k + 3 \le n$ 的非負整數 $k$ 的個數，$K = \max(0, \lfloor (n-3)/2 \rfloor
+ 1) = \lfloor (n-1)/2 \rfloor$。

$n = 2m + 1$ 時 $K = m$，$s = x_{m+1} = x_{\lceil n/2 \rceil}$；$n = 2m$ 時
$K = m - 1$，$s = x_m = x_{n/2}$。 兩種情況都是 $x_{\lceil n/2 \rceil}$。
\end{proof}

$f$ 從 $h.\mathit{next}$ 出發而不是從 $h$ 出發，是為了偶數長度時讓 $s$ 停在兩個中
間節點的\textbf{前一個}。 $n = 2$ 時這樣切成 $1 + 1$；如果 $f$ 從 $h$ 出發，會切成
$2 + 0$，合併排序用它的話會一直切不開（觀念題第 11 題）。

快慢指標和「先數長度」的做法都走了大約 $1.5n$ 個指標：$f$ 走 $n$ 步、$s$ 走 $n/2$
步。 它的好處在不需要先知道長度。 同樣的技巧也能判斷串列有沒有環：有環的話快指標會在環
上追上慢指標，沒有環的話快指標會先走到 \NIL。 第~\ref{sec:ll-cycle}~節會把這件事講完整。

\subsection{例子：兩位助教的成績名單}

下面的程式把三個操作串起來：合併兩位助教的名單，切成兩半，再把後半反轉。 節點除了
學號還有成績；為了版面只印學號，成績跟著節點一起搬，不用另外處理。 \texttt{merge} 和 \texttt{split\_half} 是虛擬
碼~\ref{alg:ll-merge}、\ref{alg:ll-split} 的直接翻譯，\texttt{merge} 的頭節點是區域
變數。

\begin{example}[label=ex:ll-ops]{c}{(list\_ops.c) 合併、切半、反轉}
#include <stdio.h>

typedef struct Node {
    int id, score;          // 學號、小考成績
    struct Node *next;
} Node;

// 把陣列 a[0..n-1] 裡的節點依序串起來，回傳第一個。
Node *link_all(Node *a, int n)
{
    for (int i = 0; i + 1 < n; i++)
        a[i].next = &a[i + 1];
    a[n - 1].next = NULL;
    return &a[0];
}

// 合併兩條依學號排好的串列，不配置新節點。*cmp 記錄比較次數。
Node *merge(Node *a, Node *b, int *cmp)
{
    Node head = {0, 0, NULL};
    Node *t = &head;
    *cmp = 0;
    while (a != NULL && b != NULL) {
        (*cmp)++;
        if (a->id <= b->id) {
            t->next = a;
            a = a->next;
        } else {
            t->next = b;
            b = b->next;
        }
        t = t->next;
    }
    t->next = (a != NULL) ? a : b;  // 剩下的一整段直接接上
    return head.next;
}

// 切成前後兩半，前半有 ceil(n/2) 個節點。回傳後半的開頭。
Node *split_half(Node *h)
{
    if (h == NULL)
        return NULL;
    Node *slow = h, *fast = h->next;
    while (fast != NULL && fast->next != NULL) {
        slow = slow->next;          // 走一步
        fast = fast->next->next;    // 走兩步
    }
    Node *back = slow->next;
    slow->next = NULL;              // slow 是前半的最後一個
    return back;
}

// 就地反轉，回傳新的開頭。
Node *reverse(Node *h)
{
    Node *done = NULL;              // 已經反轉好的一段
    while (h != NULL) {
        Node *rest = h->next;       // 先記住還沒處理的部分
        h->next = done;
        done = h;
        h = rest;
    }
    return done;
}

void print_list(const char *name, const Node *p)   // 只印學號
{
    printf("%s:", name);
    for (; p != NULL; p = p->next)
        printf(" %d", p->id);
    printf("\n");
}

int main(void)
{
    // 節點直接放在陣列裡；三個操作都只改 next，不需要 malloc。
    Node ta1[] = {{103, 72, NULL}, {108, 90, NULL}, {112, 65, NULL},
                  {120, 88, NULL}};
    Node ta2[] = {{101, 80, NULL}, {109, 77, NULL}, {110, 93, NULL},
                  {115, 58, NULL}, {117, 84, NULL}};
    Node *a = link_all(ta1, 4), *b = link_all(ta2, 5);
    print_list("TA1", a);
    print_list("TA2", b);
    int cmp;
    Node *all = merge(a, b, &cmp);
    print_list("merged", all);
    printf("comparisons = %d\n", cmp);
    Node *back = split_half(all);
    print_list("room A", all);
    print_list("room B", back);
    print_list("room B reversed", reverse(back));
    return 0;
}
\end{example}
\iosample{}{TA1: 103 108 112 120 \\ TA2: 101 109 110 115 117 \\ merged: 101 103 108 109 110 112 115 117 120 \\ comparisons = 8 \\ room A: 101 103 108 109 110 \\ room B: 112 115 117 120 \\ room B reversed: 120 117 115 112}

$m + n = 9$，比較了 $8 = m + n - 1$ 次，是最多的情況：最大的 $120$ 在助教一的名單，
第二大的 $117$ 在助教二的名單，一直到只剩 $120$ 才停。 合併後有 $9$ 個節點，\texttt{split\_half} 讓前半
有 $\lceil 9/2 \rceil = 5$ 個。 節點都放在 \texttt{main} 的兩個陣列裡，三個操作
只改 \texttt{next}，沒有配置也沒有釋放任何節點。

\section{判斷串列有沒有環}
\label{sec:ll-cycle}

前面的串列都假設最後一個節點的 $\mathit{next}$ 是 \NIL。 程式有錯時這件事不一定成
立：觀念題第 7 題 (b) 拿掉一行之後，$1$ 和 $2$ 互相指著；兩條串列接錯地方，也會讓某
個節點指回前面。 這時「走到 \NIL 為止」的迴圈永遠停不下來。 這一節的問題是：只給
$\mathit{head}$，怎麼判斷從它出發會不會走進一個環？ 有環的話，環從哪一個節點開始、
有幾個節點？ 而且希望只用 $\Theta(1)$ 的額外空間，也不改動串列。

\subsection{$\rho$ 形的串列}

從 $h$ 出發一直沿著 $\mathit{next}$ 走，依序經過的節點記成
\[
x_0 = h, \qquad x_{i+1} = x_i.\mathit{next}.
\]
沒有環時，有 $n$ 個節點的串列走 $n$ 步就到 \NIL。 有環時永遠走不完，但節點只有有限
個，總會走到一個之前走過的節點。 令 $\mu$ 是最小的足標，使得 $x_\mu$ 之後還會再出現；
再令 $\lambda$ 是最小的正整數，使得 $x_{\mu+\lambda} = x_\mu$。 $x_\mu$ 叫\textbf{環的
起點}，$\lambda$ 叫\textbf{環長}；$x_0, \ldots, x_{\mu-1}$ 是環前面的\textbf{尾巴}，
$\mu$ 是尾巴的長度（整條串列就是一個環時 $\mu = 0$）。 畫出來像希臘字母 $\rho$：先
一段直線，再一個圈。 下圖的 $\mu = 3$、$\lambda = 5$。

\begin{center}
\begin{tikzpicture}
\node[lname] (hd) at (-1.15,0) {head};
\foreach \k/\v in {0/12, 1/25, 2/31, 3/47, 4/50, 5/66, 6/73, 7/88} {
    \lnode{c\k}{\k*1.75,0}{\v}
    \node[font=\scriptsize, text=black!60] at (\k*1.75+0.2,0.55) {$x_{\k}$};
}
\draw[lnk] (hd) -- (c0);
\foreach \k [evaluate=\k as \j using int(\k+1)] in {0,...,6} {
    \llink{c\k}{-- (c\j)}
}
\fill[black!70] (c7p.center) circle (1.1pt);
\draw[lnk, rounded corners=6pt] (c7p.center) -- ++(0.5,0) -- ++(0,-0.75) -| (c3.south);
\draw[decorate, decoration={brace, amplitude=4pt}, black!60]
    (-0.45,0.85) -- node[above=3pt, font=\small] {尾巴 $\mu = 3$} (3.95,0.85);
\draw[decorate, decoration={brace, amplitude=4pt}, black!60]
    (4.8,0.85) -- node[above=3pt, font=\small] {環 $\lambda = 5$} (12.9,0.85);
\end{tikzpicture}
\diagcaption{一條有環的串列。 $88$ 的 $\mathit{next}$ 指回 $47$，所以環的起點是
$x_3 = 47$。 從 head 走下去的節點序列是 $12, 25, 31, 47, 50, 66, 73, 88, 47, 50,
\ldots$。}
\label{fig:ll-rho}
\end{center}

\begin{lemma}
\label{lem:ll-rho}
設從 $h$ 出發會走進環。 則 $x_0, \ldots, x_{\mu+\lambda-1}$ 兩兩不同，從 $h$ 走得到
的節點恰好是這 $n = \mu + \lambda$ 個。 而且對 $0 \le i < j$，
\[
x_i = x_j \iff i \ge \mu \text{ 而且 } \lambda \text{ 整除 } j - i .
\]
\end{lemma}

\begin{proof}
先證三件事。

(a) 對所有 $t \ge \mu$，$x_{t+\lambda} = x_t$。 對 $t$ 歸納：$t = \mu$ 是 $\lambda$ 的
定義；若 $x_{t+\lambda} = x_t$，兩邊各走一步得到 $x_{t+1+\lambda} = x_{t+1}$。 重複
用 (a) 就得到：$i \ge \mu$ 而且 $\lambda$ 整除 $j - i$ 時，$x_j = x_i$。

(b) 若 $x_i = x_j$、$i < j$，則 $i \ge \mu$：$x_i$ 之後又出現了，由 $\mu$ 的最小性，
$i \ge \mu$。

(c) 若 $\mu \le i < j$ 而且 $x_i = x_j$，則 $\lambda$ 整除 $j - i$。 寫成 $j - i = q\lambda
+ r$，$0 \le r < \lambda$。 由 (a)，$x_j = x_{i+r}$，所以 $x_i = x_{i+r}$。 取 $t \ge 0$
使 $i + t = \mu + s\lambda$（$s$ 是某個非負整數），兩邊各再走 $t$ 步，由 (a) 得
$x_\mu = x_{i+t} = x_{i+r+t} = x_{\mu+r}$。 如果 $r > 0$，這和 $\lambda$ 是最小的矛
盾，所以 $r = 0$。

(a)、(b)、(c) 合起來就是要證的等價式。 若 $x_0, \ldots, x_{\mu+\lambda-1}$ 裡有兩個
相同的 $x_i = x_j$（$i < j$），則 $\lambda$ 整除 $j - i$，但 $0 < j - i < \lambda$，
矛盾，所以它們兩兩不同。 由 (a)，之後的節點都在 $x_\mu, \ldots, x_{\mu+\lambda-1}$
之中重複，所以走得到的節點恰好是這 $\mu + \lambda$ 個。
\end{proof}

引理說的是：尾巴上的節點各只出現一次；環上的節點每隔 $\lambda$ 步出現一次。

\subsection{Floyd 的快慢指標}

用一個集合記住走過的節點，就能發現第一個重複的節點，但要 $\Theta(n)$ 的空間。
\textbf{Floyd 的判環法}（也叫龜兔賽跑）只用兩個指標：慢指標 $s$ 一次走一步，快指標
$f$ 一次走兩步，兩個都從 $h$ 出發。 沒有環時，$f$ 先走到 \NIL；有環時，兩個都進了環
之後，$f$ 每一輪比 $s$ 多走一步，一定會在環上追上 $s$。 追上之後，再把一個指標放回
$h$，兩個都改成一次一步，下一次碰頭的地方就是環的起點。

\begin{pseudo}[label=alg:ll-floyd]{Floyd 判環：相遇點、環的起點、環長}
\Function{Floyd-Meet}{$h$} \Comment{有環時回傳相遇的節點，沒有環回傳 \NIL}
    \State $s \gets h$；$f \gets h$
    \While{$f \ne \NIL$ \AND $f.\mathit{next} \ne \NIL$}
        \State $s \gets s.\mathit{next}$
        \State $f \gets f.\mathit{next}.\mathit{next}$
        \If{$s = f$}
            \State \Return $s$
        \EndIf
    \EndWhile
    \State \Return \NIL
\EndFunction
\Statex
\Function{Cycle-Start}{$h$}
    \State $m \gets \Call{Floyd-Meet}{h}$
    \If{$m = \NIL$}
        \State \Return \NIL \Comment{沒有環}
    \EndIf
    \State $s \gets h$；$f \gets m$ \Comment{一個從頭、一個從相遇點}
    \While{$s \ne f$}
        \State $s \gets s.\mathit{next}$；$f \gets f.\mathit{next}$ \Comment{都走一步}
    \EndWhile
    \State \Return $s$
\EndFunction
\Statex
\Function{Cycle-Length}{$c$} \Comment{$c$ 是環上的節點}
    \State $\lambda \gets 1$；$t \gets c.\mathit{next}$
    \While{$t \ne c$}
        \State $t \gets t.\mathit{next}$；$\lambda \gets \lambda + 1$
    \EndWhile
    \State \Return $\lambda$
\EndFunction
\end{pseudo}

\proc{Floyd-Meet} 的迴圈條件和 \proc{Split-Half}（虛擬碼~\ref{alg:ll-split}）一樣：
$f$ 和 $f.\mathit{next}$ 都存在，才能走兩步。 差別在兩個指標都從 $h$ 出發，而且每走
完一輪就比一次。

\begin{theorem}
\label{thm:ll-floyd}
\begin{enumerate}
\item 沒有環時（$n$ 個節點），\proc{Floyd-Meet} 的迴圈跑 $\lfloor n/2 \rfloor$ 輪後
停下，回傳 \NIL。
\item 有環時，\proc{Floyd-Meet} 一定回傳一個節點。 迴圈跑的輪數 $k^*$ 是
$\lambda$ 的倍數裡、不小於 $\max(\mu, 1)$ 的最小那一個，而且 $k^* \le n = \mu +
\lambda$。
\item 有環時，\proc{Cycle-Start} 的第二個迴圈恰好跑 $\mu$ 輪，回傳環的起點
$x_\mu$。
\end{enumerate}
\end{theorem}

\begin{proof}
先看 \proc{Floyd-Meet}：做完 $k$ 輪之後 $s = x_k$、$f = x_{2k}$。 $k = 0$ 時成立，每
一輪 $s$ 前進一格、$f$ 前進兩格，所以一直成立。

(1) 沒有環時，$x_0, \ldots, x_{n-1}$ 兩兩不同，$x_n = \NIL$。 做完 $k$ 輪後還要再
做一輪，條件是 $x_{2k}$ 和 $x_{2k+1}$ 都存在，也就是 $2k + 2 \le n$。 滿足這個條件
的 $k \ge 0$ 有 $\lfloor n/2 \rfloor$ 個，所以迴圈跑 $\lfloor n/2 \rfloor$ 輪。 中
間每次比較的是 $x_k$ 和 $x_{2k}$（$k \ge 1$），足標不同，節點也不同，第 7 行不會執行，
最後回傳 \NIL。

(2) 有環時串列上沒有 \NIL，迴圈只會從第 7 行離開。 第 $k$ 輪（$k \ge 1$）結束時
$s = f$，就是 $x_k = x_{2k}$；由引理~\ref{lem:ll-rho}（取 $i = k$、$j = 2k$），這等
價於 $k \ge \mu$ 而且 $\lambda$ 整除 $k$。 所以迴圈在第一個滿足這兩個條件的 $k \ge 1$
停下，就是定理說的 $k^*$。 $\mu = 0$ 時 $k^* = \lambda = n$。 $\mu \ge 1$ 時，
$\mu, \mu + 1, \ldots, \mu + \lambda - 1$ 這 $\lambda$ 個連續整數裡恰好有一個是
$\lambda$ 的倍數，所以 $k^* \le \mu + \lambda - 1 < n$。

(3) 令 $m = x_{k^*}$。 \proc{Cycle-Start} 的第二個迴圈做完 $i$ 輪後 $s = x_i$、
$f = x_{k^* + i}$。 $i < \mu$ 時，$x_i$ 在尾巴上，由引理~\ref{lem:ll-rho} 它和其他
足標的節點都不同，所以 $s \ne f$，迴圈繼續。 $i = \mu$ 時，兩個足標 $\mu$ 和 $k^* +
\mu$ 都 $\ge \mu$，相差 $k^*$，是 $\lambda$ 的倍數，所以 $s = f = x_\mu$，迴圈停下。
\end{proof}

證明用的是代數，下面兩段把同樣的事用「追趕」的方式說一次，考試寫說明時比較順手。

\textbf{為什麼一定會相遇，不會跳過去。} 慢指標在第 $\mu$ 輪走到環的起點，這時快指標
已經在環上。 沿著環的方向看，快指標要再追 $g$ 步才碰到慢指標（$0 \le g < \lambda$）。
之後每一輪慢指標走 $1$ 步、快指標走 $2$ 步，$g$ 恰好減少 $1$，所以 $g$ 會一路減到
$0$，中間不會從 $1$ 直接變成負的。 也就是說，再過 $g$ 輪就相遇（$g = 0$ 表示第 $\mu$
輪就相遇了；$\mu = 0$ 時兩個指標一開始就在一起，要等繞完第一圈）。

\textbf{為什麼從頭和從相遇點一起走，會在起點碰頭。} 相遇點是從 $h$ 走 $k^*$ 步到的，
$k^*$ 是 $\lambda$ 的倍數。 從相遇點再走 $\mu$ 步，等於從 $h$ 走 $k^* + \mu$ 步：先走
$\mu$ 步到起點，再繞 $k^* / \lambda$ 整圈，還是停在起點。 從 $h$ 出發的那個指標走
$\mu$ 步也剛好到起點。 兩個指標同一輪到達起點；在那之前，從頭出發的指標還在尾巴上，
另一個在環上，不可能碰頭。 注意相遇點本身\textbf{不一定}是起點：下面的例子相遇在
$66$，起點是 $47$。

\textbf{成本。} \proc{Floyd-Meet} 最多 $\lfloor n/2 \rfloor$ 或 $k^* \le n$ 輪，
\proc{Cycle-Start} 的第二個迴圈 $\mu$ 輪，\proc{Cycle-Length} 從起點出發繞一圈
$\lambda$ 輪，每輪 $\Theta(1)$，總時間 $O(n)$。 反過來，沒有環時跑
$\lfloor n/2 \rfloor$ 輪，有環時 $k^* \ge \max(\mu, \lambda) \ge n/2$，所以也是
$\Omega(n)$，合起來 $\Theta(n)$。 額外空間只有幾個指標，$\Theta(1)$，而且完全沒有
改動串列。

\textbf{手算追蹤。} 用圖~\ref{fig:ll-rho} 的串列，第 $k$ 列是 \proc{Floyd-Meet} 做完
$k$ 輪後兩個指標指著的節點：

\begin{center}
\begin{tabular}{cccc}
\toprule
輪 $k$ & $s = x_k$ & $f = x_{2k}$ & 相等？ \\
\midrule
0 & $12$ & $12$ & （還沒開始比）\\
1 & $25$ & $31$ & 否 \\
2 & $31$ & $50$ & 否 \\
3 & $47$ & $73$ & 否 \\
4 & $50$ & $47$ & 否 \\
5 & $66$ & $66$ & 相遇 \\
\bottomrule
\end{tabular}
\end{center}

$\mu = 3$、$\lambda = 5$，$5$ 的倍數裡不小於 $3$ 的最小是 $5$，所以第 $5$ 輪相遇，
和定理一致。 第 $3$ 輪慢指標剛進環時，快指標在 $73$，沿著環再走 $2$ 步才碰到慢指標
（$73 \to 88 \to 47$），所以再過 $2$ 輪相遇。 接著 \proc{Cycle-Start}：

\begin{center}
\begin{tabular}{ccc}
\toprule
第二個迴圈的輪 & $s$（從 head 出發）& $f$（從相遇點出發）\\
\midrule
0 & $12$ & $66$ \\
1 & $25$ & $73$ \\
2 & $31$ & $88$ \\
3 & $47$ & $47$ \\
\bottomrule
\end{tabular}
\end{center}

跑了 $3 = \mu$ 輪，回傳 $47$。 最後從 $47$ 繞一圈 $47 \to 50 \to 66 \to 73 \to 88 \to
47$，環長 $5$。

\begin{examnote}
判環的題目，先寫出你用的版本：兩個指標從哪裡出發、每輪先走再比還是先比再走。 有些
課本讓 $f$ 從 $h.\mathit{next}$ 出發、先比再走，這時第一階段停下時，$s$ 走了 $k$ 步，
$f$ 走了 $2k + 1$ 步，比兩倍多一步，所以第二階段要讓 $f$ 從相遇點的\textbf{下一個}節點出發，才會在起點碰頭；
照本書的寫法從相遇點出發，兩個指標會永遠碰不到（環長是 $1$ 時除外）。 追蹤時一輪一列，寫出
$s$、$f$；最後分開寫「相遇點」「環的起點」「環長」三個答案。
\end{examnote}

\subsection{和其他做法比較}

\begin{center}
\begin{tabular}{>{\raggedright\arraybackslash}p{0.34\linewidth}lll}
\toprule
做法 & 時間 & 額外空間 & 改動串列 \\
\midrule
用集合記住走過的節點（雜湊表，第~12~章）& 期望 $\Theta(n)$ & $\Theta(n)$ & 不會 \\
每個節點多一個「走過了」旗標 & $\Theta(n)$ & 每個節點一個欄位 & 會，用完要清掉 \\
Floyd 快慢指標 & $\Theta(n)$ & $\Theta(1)$ & 不會 \\
\bottomrule
\end{tabular}
\end{center}

前兩種做法一走到重複的節點就知道它是環的起點，比較直接；Floyd 的方法用多走幾輪換
到常數空間，而且串列可以是唯讀的（例如別人的資料結構，或很多程式同時在讀）。
「先數長度」行不通：有環時數不完。

同一個方法不只能用在串列上。 只要有一個「下一個是誰」的規則，例如陣列
$\mathit{next}[\,]$ 表示的串列，或亂數產生器 $x_{i+1} = (a x_i + c) \bmod M$，從
任何起點一直往下走，最後都會進入一個環，Floyd 的方法可以在不記錄整個序列的情況下
找出尾巴長度和週期。

\subsection{C 程式}

下面的程式把圖~\ref{fig:ll-rho} 的八個節點放在陣列裡，改變最後一個節點的
\texttt{next}，做出四種串列：沒有環、$\rho$ 形、整條是一個環、最後一個節點指向自
己。 每一種印出相遇點、第一階段跑了幾輪、環的起點、尾巴長度和環長。

\begin{example}[label=ex:ll-cycle]{c}{(cycle.c) Floyd 判環}
#include <stdio.h>

typedef struct Node {
    int key;
    struct Node *next;
} Node;

// 第一階段：有環時回傳相遇的節點，沒有環回傳 NULL。*rounds 記錄跑了幾輪。
Node *floyd_meet(Node *h, int *rounds)
{
    Node *slow = h, *fast = h;
    *rounds = 0;
    while (fast != NULL && fast->next != NULL) {
        slow = slow->next;          // 走一步
        fast = fast->next->next;    // 走兩步
        (*rounds)++;
        if (slow == fast)
            return slow;
    }
    return NULL;
}

// 第二階段：一個從頭、一個從相遇點，每次各走一步，碰頭的地方是環的起點。
Node *cycle_start(Node *h, Node *meet, int *tail)
{
    *tail = 0;
    while (h != meet) {
        h = h->next;
        meet = meet->next;
        (*tail)++;
    }
    return h;
}

int cycle_length(const Node *c)     // c 在環上：繞一圈數節點
{
    int len = 1;
    for (const Node *t = c->next; t != c; t = t->next)
        len++;
    return len;
}

void report(const char *name, Node *h)
{
    int rounds, tail;
    Node *meet = floyd_meet(h, &rounds);
    printf("%s: ", name);
    if (meet == NULL) {
        printf("no cycle (%d rounds)\n", rounds);
        return;
    }
    Node *start = cycle_start(h, meet, &tail);
    printf("meet %d after %d rounds, ", meet->key, rounds);
    printf("start %d, tail %d, length %d\n", start->key, tail,
           cycle_length(start));
}

int main(void)
{
    int keys[] = {12, 25, 31, 47, 50, 66, 73, 88};
    Node a[8];
    for (int i = 0; i < 8; i++) {
        a[i].key = keys[i];
        a[i].next = (i + 1 < 8) ? &a[i + 1] : NULL;
    }
    report("chain", &a[0]);
    a[7].next = &a[3];      // 88 指回 47
    report("rho", &a[0]);
    a[7].next = &a[0];      // 88 指回 12：整條是一個環
    report("ring", &a[0]);
    a[7].next = &a[7];      // 88 指向自己
    report("self", &a[0]);
    return 0;
}
\end{example}
\iosample{}{chain: no cycle (4 rounds) \\ rho: meet 66 after 5 rounds, start 47, tail 3, length 5 \\ ring: meet 12 after 8 rounds, start 12, tail 0, length 8 \\ self: meet 88 after 7 rounds, start 88, tail 7, length 1}

四行都和定理~\ref{thm:ll-floyd} 對得上：沒有環的 $8$ 個節點跑 $\lfloor 8/2 \rfloor =
4$ 輪；整條是環時 $\mu = 0$，要等 $k^* = \lambda = 8$ 輪，兩個指標都回到 $12$ 才相
遇；$88$ 指向自己時 $\mu = 7$、$\lambda = 1$，第 $7$ 輪慢指標一走到 $88$ 就和早已停
在那裡打轉的快指標相遇。 第三種情況相遇點就是起點，第二種不是，所以第二階段不能省。

\section{本章常見錯誤}
\label{sec:ll-pitfalls}

\begin{center}
\begin{tabular}{p{0.24\linewidth}p{0.50\linewidth}p{0.16\linewidth}}
\toprule
錯誤 & 例子 & 見 \\
\midrule
改指標的順序錯 & 插入時先寫 \texttt{p->next = x} 再寫 \texttt{x->next = p->next}，
$x$ 指向自己，後半段遺失 & 第~\ref{sec:ll-singly}~節 \\
釋放之後還在用 & \texttt{free(p); p = p->next;}；刪除時先 \texttt{free} 才接前後 &
第~\ref{sec:ll-singly}~節 \\
改到參數的複本 & 函式收 \texttt{Node *head}，在裡面寫 \texttt{head = ...}，呼叫端的
\texttt{head} 沒變 & 第~\ref{sec:ll-singly}~節 \\
忘了第一個節點 & 刪除或插入在最前面時沒有更新 \texttt{head} &
第~\ref{sec:ll-singly}~節 \\
環狀串列用 \NIL 當停止條件 & \texttt{while (p != NULL)} 走環，永遠停不下來 &
第~\ref{sec:ll-variants}~節 \\
雙向串列只改一個方向 & 刪除時只改 \texttt{x->prev->next}，後一個節點的
\texttt{prev} 還指著 $x$ & 第~\ref{sec:ll-variants}~節 \\
係數相消的項留在結果裡 & 多項式相加得到「$0x^7$」這種節點 &
第~\ref{sec:ll-poly}~節 \\
以為還回去的節點被清空 & 用 $\mathit{key}$ 判斷一格有沒有在用；還回去之後還有別的
指標指著它 & 第~\ref{sec:ll-avail}~節 \\
反轉時沒先記住下一個 & 先改 \texttt{h->next = d}，再寫 \texttt{h = h->next}，
回到了 $d$ & 第~\ref{sec:ll-ops}~節 \\
快慢指標的起點選錯 & 偶數長度時切成「全部＋空的」，合併排序無窮遞迴 &
第~\ref{sec:ll-ops}~節 \\
把相遇點當成環的起點 & 判環時第一階段停下的節點直接當答案；$f$ 從 $h.\mathit{next}$
出發，第二階段卻從相遇點出發 & 第~\ref{sec:ll-cycle}~節 \\
\bottomrule
\end{tabular}
\end{center}

\section{觀念題}
\label{sec:ll-exercises}

題型和研究所考試、期中期末考的筆試相近，有選擇、計算、追蹤、證明和寫程式。 解答在
本章最後。

\begin{enumerate}
\item （選擇，可能不只一個答案）一條有 $n$ 個節點的單向串列，只記
$L.\mathit{head}$，沒有頭節點、也不記最後一個節點。 下列哪些操作在最壞情況下是
$\Theta(1)$？
\begin{enumerate}
\item 在最前面插入一個節點
\item 刪除最後一個節點
\item 已知指向節點 $p$ 的指標，在 $p$ 後面插入一個節點
\item 已知指向節點 $p$ 的指標，在 $p$ 前面插入一個節點，而且串列上原有的每個節點
仍然存著原本的資料
\item 求串列的長度
\end{enumerate}

\item （追蹤）對單向串列 $L$ 執行下面的程序。
\begin{pseudo}{觀念題 2}
\Procedure{Mystery}{$L$}
    \State $p \gets L.\mathit{head}$
    \While{$p \ne \NIL$ \AND $p.\mathit{next} \ne \NIL$}
        \State $q \gets p.\mathit{next}$
        \State $p.\mathit{next} \gets q.\mathit{next}$
        \State $q.\mathit{next} \gets L.\mathit{head}$
        \State $L.\mathit{head} \gets q$
        \State $p \gets p.\mathit{next}$
    \EndWhile
\EndProcedure
\end{pseudo}
\begin{enumerate}
\item $L$ 是 $1 \to 2 \to 3 \to 4 \to 5 \to 6 \to 7$。 寫出 \textbf{while} 每一輪之後
的串列。
\item 用一句話描述這個程序對一般的串列做了什麼，並寫出 $n$ 個節點時迴圈跑幾輪。
\end{enumerate}

\item （選擇，可能不只一個答案）在雙向串列裡把節點 $x$ 接在 $p$ 後面（$p$ 後面原本
有節點），要執行下面四個指定：
\[
\text{(1)}\ x.\mathit{next} \gets p.\mathit{next} \quad
\text{(2)}\ x.\mathit{prev} \gets p \quad
\text{(3)}\ p.\mathit{next}.\mathit{prev} \gets x \quad
\text{(4)}\ p.\mathit{next} \gets x
\]
下列哪些執行順序是正確的？
(A) 1, 2, 3, 4 \quad (B) 4, 1, 2, 3 \quad (C) 2, 1, 4, 3 \quad (D) 3, 1, 2, 4 \quad
(E) 1, 4, 2, 3

\item （追蹤）用陣列模擬串列，節點池有 $6$ 格。 目前的狀態如下，「--」表示那一格的
$\mathit{key}$ 沒有意義。 $\mathit{head} = 2$，$\mathit{avail} = 3$。
\begin{center}
\begin{tabular}{ccccccc}
\toprule
編號 & 0 & 1 & 2 & 3 & 4 & 5 \\
\midrule
$\mathit{key}$ & 31 & -- & 17 & -- & -- & 8 \\
$\mathit{next}$ & 5 & 4 & 0 & 1 & $-1$ & $-1$ \\
\bottomrule
\end{tabular}
\end{center}
依序執行：刪除 $31$、在最前面插入 $44$、在最前面插入 $66$、刪除 $8$。 插入用
\proc{Get-Node}，刪除時找到前一個節點接好後用 \proc{Ret-Node} 把那一格還回去。 寫出最後的
$\mathit{key}$、$\mathit{next}$、$\mathit{head}$、$\mathit{avail}$，以及從
$\mathit{head}$ 和 $\mathit{avail}$ 出發的兩條串列。

\item （追蹤）$3 \times 4$ 的矩陣 $N$ 只有四個非零元素：$N[1][2] = 5$、
$N[2][4] = -2$、$N[3][1] = 9$、$N[3][3] = 4$。 用第~\ref{sec:ll-poly}~節的正交串列
存（每列、每行各一個頭節點，串列都是環狀的）。
\begin{enumerate}
\item 一共有幾個節點？
\item 從第 $3$ 列的頭節點沿 $\mathit{right}$ 走一圈，依序經過哪些節點？
\item 把 $N[3][3]$ 改成 $0$（刪掉那個節點），要改哪兩個指標，各改成指向誰？
\item 把 $N[2][2]$ 設成 $7$（新增節點 $y$），在第 $2$ 列和第 $2$ 行各要改哪些指標？
\end{enumerate}

\item （證明）兩條排好的串列各有 $m$、$n$ 個節點（$m, n \ge 1$），用
\proc{Merge-Lists}（虛擬碼~\ref{alg:ll-merge}）合併。 證明 $\mathit{key}$ 的比較次
數 $C$ 滿足 $\min(m, n) \le C \le m + n - 1$，並對任意的 $m$、$n$ 各給一組輸入，分
別讓 $C = \min(m, n)$ 和 $C = m + n - 1$。

\item （證明）下面是遞迴版的反轉，回傳反轉後的第一個節點。
\begin{pseudo}{觀念題 7}
\Function{Rev-Rec}{$h$}
    \If{$h = \NIL$ \OR $h.\mathit{next} = \NIL$}
        \State \Return $h$
    \EndIf
    \State $r \gets \Call{Rev-Rec}{h.\mathit{next}}$
    \State $h.\mathit{next}.\mathit{next} \gets h$
    \State $h.\mathit{next} \gets \NIL$
    \State \Return $r$
\EndFunction
\end{pseudo}
\begin{enumerate}
\item 對節點個數 $n$ 用數學歸納法證明它是對的。
\item 如果拿掉第 6 行，對 $1 \to 2 \to 3$ 執行之後會發生什麼事？
\item 它的額外空間是多少？
\end{enumerate}

\item （寫程式）用第~\ref{sec:ll-singly}~節的 \texttt{Node} 和 \texttt{Node **}
寫法，寫出
\begin{center}
\texttt{void remove\_all(Node **head, int k)}
\end{center}
刪除串列裡所有 \texttt{id} 等於 \texttt{k} 的節點並釋放它們。 要求 $O(n)$ 時間，
不能另外寫「刪第一個節點」的特殊情況。

\item （找錯）下面兩個函式各有什麼錯？ 說明怎麼改。
\begin{code}{c}{觀念題 9}
// (a) 釋放整條串列
void free_all(Node *head)
{
    for (Node *p = head; p != NULL; p = p->next)
        free(p);
}

// (b) 刪除第一個節點
void drop_first(Node *head)
{
    Node *p = head;
    head = head->next;
    free(p);
}
\end{code}

\item （設計）值日環用單向環狀串列表示，只記最後一個節點 $R.\mathit{last}$。 現在
要把它拆成兩組：從今天值日的人一直到節點 $p$（含 $p$）為第一組，$p$ 之後到
$R.\mathit{last}$ 為第二組，兩組各自成為一個環，組內的輪替順序不變。 已知 $p$ 在環上
而且 $p \ne R.\mathit{last}$。 寫出 $\Theta(1)$ 的程序 \proc{Ring-Split}$(R, p)$，
回傳兩個環各自的最後一個節點。 如果環改成只記第一個節點，還能 $\Theta(1)$ 做到嗎？

\item （計算）把 \proc{Split-Half}（虛擬碼~\ref{alg:ll-split}）第 4 行的
$f \gets h.\mathit{next}$ 改成 $f \gets h$，其他不變。
\begin{enumerate}
\item $n = 1, 2, \ldots, 6$ 時，迴圈結束時 $s$ 是第幾個節點？ 寫出一般的公式。
\item 用改過的版本切半，再拿去做合併排序（切成兩半、各自遞迴排序、再合併），$n = 2$
時會發生什麼事？
\end{enumerate}

\item （計算）一條串列有 $10$ 個節點 $x_0, x_1, \ldots, x_9$，$x_0$ 是 head，
$x_i.\mathit{next} = x_{i+1}$（$i \le 8$），$x_9.\mathit{next} = x_6$。
\begin{enumerate}
\item 寫出這條串列的 $\mu$、$\lambda$。 \proc{Floyd-Meet}（虛擬碼~\ref{alg:ll-floyd}）
跑幾輪、在哪個節點相遇？ 列出每一輪之後的 $s$、$f$。
\item \proc{Cycle-Start} 的第二個迴圈跑幾輪？
\item 把 $x_9.\mathit{next}$ 改成 $x_0$，(a)、(b) 的答案變成什麼？
\item 節點數固定是 $10$，$x_9.\mathit{next}$ 可以指向 $x_0, \ldots, x_9$ 任何一個。
哪一種讓 \proc{Floyd-Meet} 跑最多輪、哪一種最少？ 各幾輪？
\end{enumerate}

\item （設計）兩條沒有環的單向串列，開頭分別是 $a$、$b$，長度 $m$、$n$。 它們可能從
某個節點 $X$ 開始共用後半段（$X$ 以後的節點兩條串列都有，兩條分支在 $X$ 合成一條，像英文字母 Y）。
寫出一個 $O(m + n)$ 時間、$O(1)$ 額外空間的程序，回傳第一個共用的節點 $X$，沒有共
用的節點時回傳 \NIL。 可以暫時改動指標，但結束前要改回來。 （提示：用第~\ref{sec:ll-cycle}~節的
\proc{Cycle-Start}。）
\end{enumerate}

\section{歷屆試題}
\label{sec:ll-pastexam}

以下是研究所入學考試的原題，題目照原卷（語言、記號、配分照錄）。 詳解在本章最後。

\begin{pastexam}[label={px:ll-1}]{臺大資工所 110 學年度 第 VII 題}
（原卷為單選題，每題答對 5 分，答錯倒扣 2.5 分。）

There is a linked-list with $n$ elements, which are $x_1,x_2,\ldots,x_n$. Let $f(n)$ be
the time complexity of inserting a new element $x_{n+1}$ to the tail of the list, and
$g(n)$ be that of deleting the last element $x_n$. If there is only one pointer pointed
to the head of the list, we know to access the list is time-consuming, and functions
$f(n)$ and $g(n)$ are denoted as $f_1(n)$ and $g_1(n)$ in this case. To improve it, we
may also add another pointer pointed to the tail of it, so the functions $f(n)$ and
$g(n)$ become $f_2(n)$ and $g_2(n)$. Which of the following option is false if $n$
becomes large. \underline{\makebox[5em]{(13)}}

(A) $\frac{f_1(n)}{g_1(n)}=1$ \qquad (B) $\frac{f_2(n)}{g_2(n)}=n$ \qquad
(C) $\frac{f_1(n)}{f_2(n)}=n$ \qquad (D) $\frac{g_1(n)}{g_2(n)}=1$
\end{pastexam}

\begin{pastexam}[label={px:ll-2}]{成大資訊聯招 114 學年度 第 2 題}
{[3\%]} Given a node structure as shown below:
\begin{minted}[fontsize=\small]{c}
typedef struct Node {
    int data;
    struct Node* next;
} Node;
\end{minted}
What will the following function do when used on a singly linked list?
\begin{minted}[fontsize=\small]{c}
void unknownFunction(Node* node) {
    if (node == NULL || node->next == NULL) {
        return;
    }
    node->data = node->next->data;
    Node* temp = node->next;
    node->next = node->next->next;
    free(temp);
}
\end{minted}
\begin{enumerate}[label=(\Alph*)]
\item The function will delete the specified node, including the last node in the list.
\item The function will delete the specified node only if it is a middle node or any node
except the last one.
\item The function will produce a segmentation fault when applied to the last node in
the list because it attempts to access node->next->data.
\item The function will remove the node but corrupt the list structure if the deleted
node has duplicate values in the list.
\item None of the above is correct.
\end{enumerate}
\end{pastexam}

\begin{pastexam}[label={px:ll-3}]{中央資工所 113 學年度 第 8 題}
（原卷為複選題，每題 5 分，全部答對才給分，答錯倒扣一分。）

Consider the following function that takes reference to head of a Doubly Linked List as
parameter. Assume that a node of doubly linked list has previous pointer as prev and next
pointer as next.
\begin{minted}[fontsize=\small]{c}
void fun(struct node **head_ref)
{
    struct node *temp = NULL;
    struct node *current = *head_ref;

    while (current != NULL)
    {
        temp = current->prev;
        current->prev = current->next;
        current->next = temp;
        current = current->prev;
    }

    if(temp != NULL )
        *head_ref = temp->prev;
}
\end{minted}
Assume that reference of head of a doubly linked list L is passed to above function. What
is the resulting linked list after the function call?
\begin{enumerate}[label=(\Alph*)]
\item \texttt{2 <--> 1 <--> 4 <--> 3 <--> 6 <-->5} when L is
\texttt{2 <--> 1 <--> 4 <--> 3 <--> 6 <-->5}
\item \texttt{5 <--> 4 <--> 3 <--> 2 <--> 1 <-->6} when L is
\texttt{5 <-->4 <--> 3 <-->2 <--> 6<-->1}
\item \texttt{6 <--> 5 <--> 4 <--> 3 <--> 2 <--> 1} when L is
\texttt{1 <--> 2 <--> 3 <--> 4 <--> 5 <-->6}
\item \texttt{6 <--> 5 <--> 4 <--> 3 <--> 1 <--> 2} when L is
\texttt{1 <--> 2 <--> 3 <--> 4 <--> 5 <-->6}
\end{enumerate}
\end{pastexam}

\begin{pastexam}[label={px:ll-4}]{臺科大資工所 112 學年度 第 2 題}
(10\%) Consider the function \textbf{func} as shown below, where the two lists list1 and
list2 are both ordered in ascending order.
\begin{minted}[fontsize=\small]{c}
struct item
{
   int value;
   struct item *next;
};

int func(struct item *list1, struct item *list2)
{
   struct item *ptr1, *ptr2;
   int count = 0;
   ptr1 = list1;
   ptr2 = list2;
   while ((ptr1 != NULL)&&(ptr2 != NULL))
   {
      if (ptr1->value < ptr2->value)
         ptr1 = ptr1 ->next;
      else if (ptr1->value > ptr2->value)
         ptr2 = ptr2->next;
      else
      {
         ptr1 = ptr1->next;
         ptr2 = ptr2->next;
         count++;
      }
   }
   return count;
}
\end{minted}
If $m$ and $n$ are the number of items in list1 and list2 respectively, what is the
worst-case time complexity of func(list1, list2)?
\end{pastexam}

\begin{pastexam}[label={px:ll-5}]{政大資科所 108 學年度 第 9 題}
(15\%) Design an algorithm for the following task.

\quad Input: a linked list.\\
\quad Task: determine whether the linked list contains a cycle.\\
\quad For example, the following linked list contains a cycle.
\begin{center}
\begin{tikzpicture}[>={Stealth[length=2.2mm]},thick,b/.style={draw,minimum size=6mm}]
\foreach \i in {1,...,5} \node[b] (n\i) at (1.2*\i,0) {};
\foreach \i/\j in {1/2,2/3,3/4,4/5} \draw[->] (n\i) -- (n\j);
\draw[->] (n5.south) to[out=-130,in=-50] (n1.south);
\end{tikzpicture}
\end{center}
\end{pastexam}

\section{實作練習}
\label{sec:ll-practice}

下面是 Aki-Judge 上和本章相關的題目。 大部分題目只要交函式和標頭檔，不寫
\texttt{main}；評測程式會檢查你有沒有真的用節點和指標，有幾題還規定不能配置新節點，
只能重接原本的 \texttt{next}。

\begin{center}
\begin{tabular}{llc}
\toprule
小節 & 題目 & 難度 \\
\midrule
單向串列 & \akiref{build-and-print-a-list}{建立並列印串列} & 暖身 \\
單向串列 & \akiref{insert-and-delete-at-a-position}{依位置插入與刪除} & 基本 \\
雙向串列 & \akiref{doubly-linked-list}{雙向鏈結串列} & 基本 \\
雙向串列 & \akiref{move-to-front-in-a-doubly-linked-list}{雙向鏈結串列的移到最前端} & 基本 \\
環狀串列 & \akiref{cutting-a-bead-ring-by-color}{依顏色拆開珠串環} & 進階 \\
用串列表示資料 & \akiref{bead-runs-pack-and-locate}{串珠壓縮與定位} & 基本 \\
可用空間 & \akiref{best-fit-memory-pool}{最佳適配記憶體池} & 進階 \\
反轉 & \akiref{reverse-a-linked-list-in-place-iterative}{就地反轉鏈結串列（迴圈）} & 暖身 \\
反轉 & \akiref{reverse-a-linked-list-in-place-recursive}{遞迴反轉鏈結串列} & 基本 \\
合併 & \akiref{merge-two-sorted-linked-lists}{合併兩條已排序的鏈結串列} & 基本 \\
合併 & \akiref{merge-chains-drop-shared-values}{合併兩條鏈並剔除共同值} & 進階 \\
重接指標 & \akiref{riffle-two-chains-by-counts}{依張數交錯洗牌兩條鏈} & 進階 \\
重接指標 & \akiref{follow-the-card}{照卡片指示抽牌} & 進階 \\
重接指標 & \akiref{threading-replies-under-posts}{把回覆接到貼文底下} & 進階 \\
綜合 & \akiref{clean-up-a-sign-up-list}{整理報名名單} & 進階 \\
\bottomrule
\end{tabular}
\end{center}

「依張數交錯洗牌兩條鏈」和「合併兩條鏈並剔除共同值」的鏈是用指標陣列表示的，
\texttt{next} 欄位換成了另一個陣列，但想法和本章一樣。 「整理報名名單」的資料量
大，要在串列上做 $O(n \lg n)$ 的排序，可以先讀第~14~章串列版的合併排序再回來寫。

\section{讀完你應該會}
\label{sec:ll-checklist}

\begin{itemize}
\item 能畫出單向、環狀、雙向串列插入和刪除時每一步的指標變化，並說出各種操作在不同
串列上的成本，以及和陣列比較的取捨。
\item 能用頭節點或 \texttt{Node **} 寫出沒有特殊情況的插入與刪除，並且釋放所有節
點，通過 AddressSanitizer 的檢查。
\item 能用串列做多項式相加並說明比較次數的上限，能畫出稀疏矩陣的正交串列並說出它
的空間和操作成本。
\item 能用陣列模擬串列和可用空間串列，逐步追蹤 $\mathit{head}$、$\mathit{avail}$ 和
$\mathit{next}[\,]$。
\item 能寫出反轉、合併、快慢指標切半，並用迴圈不變量證明反轉的正確性、算出切半時
前半的長度。
\item 能用 Floyd 的快慢指標判斷串列有沒有環、找出環的起點和環長，說明快慢指標為什麼一
定相遇、為什麼從頭和從相遇點一起走會在起點碰頭，並算出每個階段跑幾輪。
\end{itemize}

\textbf{延伸閱讀}（書末會整理完整的參考資料）

\begin{itemize}
\item Thomas H. Cormen 等，\textit{Introduction to Algorithms}，第 4 版，第 10.2 節：
雙向串列的搜尋、插入、刪除，以及用哨兵（sentinel）去掉邊界情況，對應本章的頭節點。
\item Donald E. Knuth，\textit{The Art of Computer Programming}，第 1 冊，第 2.2.3 到
2.2.6 節：鏈結配置、環狀串列、雙向串列和正交串列的原始講法，可用空間串列也在這
裡。 寫法比較老，適合想看這些結構來源的人。
\item Donald E. Knuth，\textit{The Art of Computer Programming}，第 2 冊，第 3.1 節習題
6：Floyd 判環法最早的出處，放在「亂數產生器的週期」的脈絡下，正好是第~\ref{sec:ll-cycle}~節最後說的推廣。
\item Steven S. Skiena，\textit{The Algorithm Design Manual}，第 3 版，第 3.1 節：
連續存放和鏈結存放的比較，偏重實務上的取捨。
\item 《C 程式設計基礎》第~19~章：串列的不變式、頭端插入、指標的指標、串接和有序
合併，C 語法的部分寫得比本章詳細。
\end{itemize}

\addcontentsline{toc}{section}{觀念題解答}
\section*{觀念題解答}

\begin{enumerate}
\item (a)、(c)。 (a) 和 (c) 都只改兩個指標。 (b) 要把倒數第二個節點的
$\mathit{next}$ 設成 \NIL，只能從頭走過去找，$\Theta(n)$。 (d) 要改 $p$ 前一個節點
的 $\mathit{next}$，一樣要從頭找。 （如果允許搬資料，可以在 $p$ 後面插入新節點，
再把 $p$ 和新節點的資料對調，$\Theta(1)$；但這樣 $p$ 這個節點存的資料變了，其他指
著 $p$ 的指標會看到不同的東西，題目已經排除這種做法。） (e) 要數過每個節點，
$\Theta(n)$。

\item
\begin{enumerate}
\item 每一輪把 $p$ 後面那個節點 $q$ 拿下來放到最前面，$p$ 再往前一格。
\begin{center}
\begin{tabular}{clc}
\toprule
輪 & 這一輪之後的串列 & 下一輪的 $p$ \\
\midrule
1 & $2 \to 1 \to 3 \to 4 \to 5 \to 6 \to 7$ & 3 \\
2 & $4 \to 2 \to 1 \to 3 \to 5 \to 6 \to 7$ & 5 \\
3 & $6 \to 4 \to 2 \to 1 \to 3 \to 5 \to 7$ & 7 \\
\bottomrule
\end{tabular}
\end{center}
第 3 輪之後 $p = 7$，$p.\mathit{next} = \NIL$，迴圈停。
\item 原本在偶數位置的節點，依相反的順序搬到最前面；奇數位置的節點保持原本的順序
接在後面。 $p$ 依序是原本的第 $1, 3, 5, \ldots$ 個節點，每輪處理一個偶數位置，迴
圈跑 $\lfloor n/2 \rfloor$ 輪，總時間 $\Theta(n)$。
\end{enumerate}

\item (A)、(D)。 (3) 和 (1) 都讀了 $p.\mathit{next}$，要的是原本的後一個節點，所以
都要在 (4) 之前；(2) 放在哪裡都可以。 (B)：(4) 先做，(1) 讓 $x.\mathit{next} = x$。
(C)：(3) 在 (4) 之後，$p.\mathit{next}$ 已經是 $x$，結果是 $x.\mathit{prev} = x$，
原本後一個節點的 $\mathit{prev}$ 沒有改。 (E)：同 (C) 的問題。

\item 四步之後：
\begin{center}
\begin{tabular}{ccccccc}
\toprule
編號 & 0 & 1 & 2 & 3 & 4 & 5 \\
\midrule
$\mathit{key}$ & 44 & -- & 17 & 66 & -- & 8 \\
$\mathit{next}$ & 2 & 4 & $-1$ & 0 & $-1$ & 1 \\
\bottomrule
\end{tabular}
\end{center}
$\mathit{head} = 3$，$\mathit{avail} = 5$。 串列：$66, 44, 17$（格 $3, 0, 2$）；可用：
格 $5, 1, 4$。 過程：刪除 $31$（格 $0$）時 $\mathit{next}[2] \gets 5$，格 $0$ 還回去，
$\mathit{avail} = 0 \to 3 \to 1 \to 4$；插入 $44$ 拿到格 $0$；插入 $66$ 拿到格 $3$；
刪除 $8$（格 $5$）時 $\mathit{next}[2] \gets -1$，格 $5$ 放到可用空間串列的最前面。
格 $5$ 的 $\mathit{key}$ 還是 $8$，但它已經不在串列上。

\item
\begin{enumerate}
\item $4$ 個元素節點，加上 $3$ 個列頭、$4$ 個行頭，共 $11$ 個。
\item 列頭 R3 $\to (3,1)\,9 \to (3,3)\,4 \to$ R3。
\item $(3,3)$ 是第 $3$ 列的最後一個，也是第 $3$ 行唯一的元素。 第 $3$ 列：
$(3,1).\mathit{right} \gets \text{R3}$。 第 $3$ 行：$\text{C3}.\mathit{down} \gets
\text{C3}$（第 $3$ 行變成空的）。
\item 第 $2$ 列原本只有 $(2,4)$，$y$ 的行號 $2 < 4$，要放在它前面：
$y.\mathit{right} \gets (2,4)$，$\text{R2}.\mathit{right} \gets y$。 第 $2$ 行原本只
有 $(1,2)$，$y$ 的列號 $2 > 1$，放在它後面：$y.\mathit{down} \gets \text{C2}$，
$(1,2).\mathit{down} \gets y$。
\end{enumerate}

\item 每次比較之後，恰好有一個節點接到結果上，而且以後不再參與比較。 迴圈在 $a$ 或
$b$ 變成 \NIL 時停。

下界：迴圈停下時，某一條串列的節點已經全部經由比較接上了，所以 $C \ge \min(m, n)$。

上界：迴圈停下時，另一條串列至少還剩一個節點沒有比較就整段接上，經由比較接上的節點
最多 $m + n - 1$ 個，所以 $C \le m + n - 1$。

達到下界：$a$ 是 $1, 2, \ldots, m$，$b$ 是 $m + 1, \ldots, m + n$，而且 $m \le n$。
每次比較都接上 $a$ 的節點，$m$ 次後 $a$ 用完，$C = m = \min(m, n)$（$m > n$ 時把兩條
對調）。

達到上界：讓兩條串列的最大值分別是全部的最大和第二大，例如 $a$ 是
$1, 2, \ldots, m - 1, m + n$，$b$ 是 $m, m + 1, \ldots, m + n - 1$。 $b$ 的最後一個
節點 $m + n - 1$ 接上之前，$a$ 一直還有 $m + n$ 沒接，所以前 $m + n - 1$ 個節點都是
經由比較接上的，$C = m + n - 1$。

\item
\begin{enumerate}
\item 要證明的敘述：對 $n$ 個節點的串列 $x_1 \to \cdots \to x_n$（$n = 0$ 時 $h =
\NIL$），\proc{Rev-Rec}$(h)$ 回傳 $x_n$，執行完之後串列是 $x_n \to \cdots \to x_1 \to
\NIL$。

基底：$n = 0$ 或 $n = 1$，第 2、3 行直接回傳 $h$，空串列和一個節點的串列反轉後不變。

歸納：設 $n \ge 2$，而且敘述對 $n - 1$ 個節點成立。 第 4 行對 $x_2 \to \cdots \to
x_n$ 遞迴，由歸納假設，回傳 $r = x_n$，而且 $x_n \to \cdots \to x_2 \to \NIL$。 遞迴
只改了 $x_2, \ldots, x_n$ 的 $\mathit{next}$，所以 $x_1.\mathit{next}$ 仍然是
$x_2$。 第 5 行讓 $x_2.\mathit{next} = x_1$，第 6 行讓 $x_1.\mathit{next} = \NIL$，
串列變成 $x_n \to \cdots \to x_2 \to x_1 \to \NIL$，回傳 $x_n$。

\item 這裡第 6 行指的是 $h.\mathit{next} \gets \NIL$。 拿掉之後，最外層呼叫結束時
$x_1.\mathit{next}$ 仍然是 $x_2$，而 $x_2.\mathit{next} = x_1$，串列變成 $3 \to 2
\to 1 \to 2 \to 1 \to \cdots$，$1$ 和 $2$ 互相指著，形成一個環。 走訪它的程式會停不
下來。 （內層呼叫留下的 $x_i.\mathit{next} = x_{i+1}$ 不會出問題，因為上一層的第 5
行會把它改成 $x_{i-1}$；只有最外層的 $x_1$ 沒有上一層來改。）

\item $n \ge 1$ 時呼叫最多同時疊了 $n$ 層（$x_1$ 到 $x_n$ 各一層），每層 $\Theta(1)$，
額外空間 $\Theta(n)$；迴圈版只要 $\Theta(1)$。
\end{enumerate}

\item 用 \texttt{link} 指著「目前要檢查的節點的那個指標欄位」。 刪掉時
\texttt{link} 不動（它現在指著下一個節點），沒刪才往前一格。
\begin{code}{c}{觀念題 8 的解答}
void remove_all(Node **head, int k)
{
    Node **link = head;
    while (*link != NULL) {
        if ((*link)->id == k) {
            Node *dead = *link;
            *link = dead->next;     // 跳過 dead
            free(dead);
        } else {
            link = &(*link)->next;
        }
    }
}
\end{code}
每個節點檢查一次，$O(n)$。 第一個節點要刪時，\texttt{link} 就是 \texttt{head}，
\texttt{*link = dead->next} 直接更新了呼叫端的 \texttt{head}，不需要特殊情況。 連續
好幾個節點都要刪時，\texttt{link} 一直不動，也是對的。

\item (a) \texttt{free(p)} 之後，迴圈的 \texttt{p = p->next} 讀了已經釋放的記憶體。
要先記住下一個再釋放：迴圈裡寫 \texttt{Node *next = p->next; free(p); p = next;}，
和範例~\ref{ex:ll-poly} 的 \texttt{print\_and\_free} 一樣。

(b) 兩個錯。 第一，\texttt{head} 是參數的複本，改它不會改到呼叫端；函式回來之後，
呼叫端的 \texttt{head} 還指著已經釋放的節點。 第二，空串列時 \texttt{head->next}
讀了 \NIL 的欄位。 改成收 \texttt{Node **head}：先檢查 \texttt{*head != NULL}，再
寫 \texttt{Node *p = *head; *head = p->next; free(p);}，呼叫時傳
\texttt{\&head}。 另一種改法是回傳新的開頭，呼叫端寫
\texttt{head = drop\_first(head);}。

\item 第一組的第一個人是原本的 $f = R.\mathit{last}.\mathit{next}$，最後一個人是 $p$；
第二組的第一個人是 $g = p.\mathit{next}$，最後一個人是原本的 $R.\mathit{last}$。 拆開
之後，$p$ 要接回 $f$，$R.\mathit{last}$ 要接到 $g$。
\begin{pseudo}{觀念題 10 的解答}
\Function{Ring-Split}{$R, p$}
    \State $f \gets R.\mathit{last}.\mathit{next}$ \Comment{今天值日的人}
    \State $g \gets p.\mathit{next}$ \Comment{第二組的第一個人}
    \State $p.\mathit{next} \gets f$
    \State $R.\mathit{last}.\mathit{next} \gets g$
    \State \Return $(p, R.\mathit{last})$
\EndFunction
\end{pseudo}
只讀兩個、改兩個指標，$\Theta(1)$。 $p$ 就是 $f$ 的時候（第一組只有一個人），第 4 行讓
$p$ 指向自己，仍然正確。 只記第一個節點時，第 5 行要改的是最後一個節點的
$\mathit{next}$，得繞一圈才找得到，變成 $\Theta(n)$。

\item
\begin{enumerate}
\item 做完 $k$ 輪後 $s = x_{1+k}$、$f = x_{1+2k}$。 迴圈繼續的條件是 $x_{1+2k}$ 和
$x_{2+2k}$ 都存在，即 $2k + 2 \le n$，所以跑 $\lfloor n/2 \rfloor$ 輪，結束時
$s = x_{\lfloor n/2 \rfloor + 1}$。 $n = 1, \ldots, 6$ 依序是第 $1, 2, 2, 3, 3, 4$
個。 $n$ 是奇數時和原來一樣停在正中間；$n$ 是偶數時停在兩個中間節點的後一個。
\item $n = 2$ 時 $s = x_2$，在 $s$ 後面切開，前半是 $x_1, x_2$，後半是空的。 合併排序
對前半遞迴，又是同一條長度 $2$ 的串列，永遠切不小，無窮遞迴（實際執行會用光呼叫堆
疊）。 所以給合併排序用的切法，偶數長度時一定要停在前一個中間節點，也就是原本的
$f \gets h.\mathit{next}$。
\end{enumerate}

\item
\begin{enumerate}
\item 尾巴是 $x_0, \ldots, x_5$，$\mu = 6$；環是 $x_6 \to x_7 \to x_8 \to x_9 \to x_6$，
$\lambda = 4$。 做完 $k$ 輪後 $s = x_k$、$f = x_{2k}$；足標 $t \ge 6$ 時，$x_t$ 就是
$x_{6 + ((t - 6) \bmod 4)}$。
\begin{center}
\begin{tabular}{ccccccccc}
\toprule
輪 $k$ & 1 & 2 & 3 & 4 & 5 & 6 & 7 & 8 \\
\midrule
$s$ & $x_1$ & $x_2$ & $x_3$ & $x_4$ & $x_5$ & $x_6$ & $x_7$ & $x_8$ \\
$f$ & $x_2$ & $x_4$ & $x_6$ & $x_8$ & $x_6$ & $x_8$ & $x_6$ & $x_8$ \\
\bottomrule
\end{tabular}
\end{center}
第 $8$ 輪相遇在 $x_8$。 用定理~\ref{thm:ll-floyd}：$4$ 的倍數裡不小於 $6$ 的最小是
$8$，一致。
\item $\mu = 6$ 輪：$s$ 從 $x_0$、$f$ 從 $x_8$ 出發，$6$ 輪後 $s = x_6$，$f = x_{14} =
x_6$。
\item $\mu = 0$、$\lambda = 10$。 \proc{Floyd-Meet} 要等 $k^* = 10$ 輪，相遇在
$x_{10} = x_0$；第二個迴圈一開始 $s = f = x_0$，跑 $0$ 輪，起點是 $x_0$。
\item 指向 $x_\mu$ 時尾巴長 $\mu$、環長 $10 - \mu$，$k^*$ 是 $10 - \mu$ 的倍數裡不小於
$\max(\mu, 1)$ 的最小那一個：
\begin{center}
\begin{tabular}{ccccccccccc}
\toprule
$\mu$ & 0 & 1 & 2 & 3 & 4 & 5 & 6 & 7 & 8 & 9 \\
\midrule
$k^*$ & 10 & 9 & 8 & 7 & 6 & 5 & 8 & 9 & 8 & 9 \\
\bottomrule
\end{tabular}
\end{center}
最多 $10$ 輪（指向 $x_0$，整條是一個環），最少 $5$ 輪（指向 $x_5$）。 一般來說
$k^*$ 是 $\lambda$ 的正倍數，又不小於 $\mu$，所以 $k^* \ge \max(\mu, \lambda) \ge
\lceil n/2 \rceil$；$\mu = \lfloor n/2 \rfloor$ 時 $\lambda = \lceil n/2 \rceil \ge \mu$，
$k^* = \lambda$，下界達到。 上界 $k^* \le n$ 只在 $\mu = 0$ 時達到。
\end{enumerate}

\item 想法：把 $a$ 那條串列接成一個環，再從 $b$ 出發找環的起點。
\begin{pseudo}{觀念題 13 的解答}
\Function{First-Common}{$a, b$}
    \If{$a = \NIL$ \OR $b = \NIL$}
        \State \Return \NIL
    \EndIf
    \State $t \gets a$
    \While{$t.\mathit{next} \ne \NIL$} \Comment{找 $a$ 的最後一個節點}
        \State $t \gets t.\mathit{next}$
    \EndWhile
    \State $t.\mathit{next} \gets a$ \Comment{暫時接成環}
    \State $X \gets \Call{Cycle-Start}{b}$
    \State $t.\mathit{next} \gets \NIL$ \Comment{改回來}
    \State \Return $X$
\EndFunction
\end{pseudo}
正確性：接成環之後，環上的節點恰好是 $a$ 那條串列的全部節點。 如果兩條串列從 $X$
開始共用，從 $b$ 出發會先經過 $b$ 自己的節點（不在 $a$ 上，也就不在環上），第一個
走到的環上節點就是 $X$；之後一直在環上打轉。 所以從 $b$ 看，這是一條尾巴是 $b$ 獨有
的那一段、環的起點是 $X$ 的 $\rho$ 形串列，由定理~\ref{thm:ll-floyd}，
\proc{Cycle-Start} 回傳 $X$。 如果沒有共用的節點，從 $b$ 出發永遠碰不到環，走到
\NIL，\proc{Cycle-Start} 回傳 \NIL。 第 9 行在回傳之前把 $t.\mathit{next}$ 改回
\NIL，兩條串列恢復原狀。

時間：找 $t$ 走 $m - 1$ 步；從 $b$ 走得到的節點最多 $m + n$ 個，\proc{Cycle-Start}
是 $O(m + n)$。 合計 $O(m + n)$，額外空間只有幾個指標。

不改動指標的做法：先各走一次求出 $m$、$n$ 和兩條的最後一個節點；最後一個節點不同
就沒有共用。 相同的話，讓長的那條先走 $|m - n|$ 步，再兩條一起一次一步，第一次指到
同一個節點就是 $X$（共用的後半段等長，所以兩個指標會同時到 $X$）。 一樣是
$O(m + n)$ 時間、$O(1)$ 空間。
\end{enumerate}

\addcontentsline{toc}{section}{歷屆試題詳解}
\section*{歷屆試題詳解}

\pastsol{px:ll-1}{臺大資工所 110 學年度 第 VII 題}
答案是 (B)。

題目沒有說是哪一種串列，但它說刪掉最後一個元素要花時間、而且只多一個 tail 指標，
這裡照單向串列來解（最後說明為什麼不是雙向）。 選項寫的「$= 1$」「$= n$」不能照字面
當成比值，應該讀成成長的階數：「$= 1$」是分子和分母同階，「$= n$」是分子比分母多一
個 $n$ 的因子。

四個函數（第~\ref{sec:ll-variants}~節的比較表就是這四格）：
\begin{itemize}
\item $f_1(n) = \Theta(n)$：只有 head，要把 $x_{n+1}$ 接到 $x_n$ 後面，得先從 head 走
$n - 1$ 步找到 $x_n$。
\item $g_1(n) = \Theta(n)$：刪 $x_n$ 要把 $x_{n-1}.\mathit{next}$ 設成 \NIL，得先從
head 走到 $x_{n-1}$。
\item $f_2(n) = \Theta(1)$：$\mathit{tail}.\mathit{next} \gets x_{n+1}$，
$\mathit{tail} \gets x_{n+1}$。
\item $g_2(n) = \Theta(n)$：刪掉 $x_n$ 之後 tail 要改成指向 $x_{n-1}$，單向串列不能從
$x_n$ 往回走，還是得從 head 找 $x_{n-1}$。 多一個 tail 指標對刪除最後一個沒有幫助。
\end{itemize}

逐項檢查：
\begin{itemize}
\item (A) $f_1 / g_1 = \Theta(n) / \Theta(n)$，同階，對。
\item (B) $f_2 / g_2 = \Theta(1) / \Theta(n) = \Theta(1/n)$，不是 $n$，\textbf{錯}。 加了 tail 之後
變快的只有插入，刪除沒有變。
\item (C) $f_1 / f_2 = \Theta(n) / \Theta(1) = \Theta(n)$，對。
\item (D) $g_1 / g_2 = \Theta(n) / \Theta(n)$，同階，對。
\end{itemize}

如果是雙向串列，$g_2$ 會變成 $\Theta(1)$（$x_n.\mathit{prev}$ 就是 $x_{n-1}$），這
時 (D) 的比值是 $n$，也錯了，單選題會有兩個答案。 所以題意是單向串列。

\pastsol{px:ll-2}{成大資訊聯招 114 學年度 第 2 題}
答案是 (B)。

函式做的事：把下一個節點的資料抄到 \texttt{node} 裡，再把下一個節點從串列上拿掉、
釋放。 被釋放的其實是下一個節點，但串列上少掉的是 \texttt{node} 原本的那筆資料，看
起來就像把 \texttt{node} 刪掉了。 觀念題第 1 題 (d) 的解答用同一招在 $p$ 前面插入，
這裡反過來用在刪除。 例如串列 $3 \to 8 \to 5$：
\begin{itemize}
\item 對存 $8$ 的節點呼叫：它的 \texttt{data} 變成 $5$，\texttt{next} 改成原本 $5$
那個節點的 \texttt{next}（\texttt{NULL}），再釋放原本存 $5$ 的節點。 結果是 $3 \to 5$，
正是「刪掉 $8$」。
\item 對最後一個節點（存 $5$）呼叫：\texttt{node->next == NULL}，直接返回，串列不變。
\end{itemize}

逐項檢查：
\begin{itemize}
\item (A) 錯。 對最後一個節點，函式第一個 \texttt{if} 就返回了，什麼都沒刪。 這也不
是寫法的疏忽：刪掉最後一個節點要把它前一個節點的 \texttt{next} 設成 \texttt{NULL}，
函式只拿到 \texttt{node}，沒辦法找到前一個。
\item (B) 對。 只要 \texttt{node} 不是最後一個，串列上就少了它的資料。
\item (C) 錯。 函式一開始就檢查 \texttt{node->next == NULL}，最後一個節點在這裡就
返回了，後面的 \texttt{node->next->data} 根本沒有執行，不會 segmentation fault。
\item (D) 錯。 函式沒有用資料值去找任何東西，串列裡有沒有重複的值都一樣；它改的兩
個地方（\texttt{node->data}、\texttt{node->next}）都接得好好的，結構沒有壞。
\item (E) 錯，因為 (B) 對。
\end{itemize}

補充：這個技巧在考卷上常出現，但實務上要小心。 真正被 \texttt{free} 的是下一個節點，
如果程式別的地方還留著指向它的指標（例如串列另外記了 tail，而 \texttt{node} 正好是
倒數第二個），那個指標就指到已經釋放的記憶體；指向 \texttt{node} 的指標也會發現它
存的資料變了。

\pastsol{px:ll-3}{中央資工所 113 學年度 第 8 題}
答案是 (C)。

\texttt{fun} 對每個節點把 \texttt{prev}、\texttt{next} 兩個指標對調，然後讓
\texttt{current} 走到對調後的 \texttt{prev}，也就是原本的下一個節點，所以每個節點都
會被處理一次。 每個節點兩個方向都反過來，整條串列就反轉了。 迴圈結束時
\texttt{temp} 是最後處理的節點 $x_n$ 原本的 \texttt{prev}，也就是 $x_{n-1}$；
$x_{n-1}$ 對調之後 \texttt{prev} 是原本的 \texttt{next}，也就是 $x_n$，所以
\texttt{*head\_ref = temp->prev} 讓新的開頭是 $x_n$。

用 L $= 1 \leftrightarrow 2 \leftrightarrow 3 \leftrightarrow 4 \leftrightarrow 5
\leftrightarrow 6$ 追蹤，每一列是迴圈的一輪（「之後」是對調後的值）：
\begin{center}
\begin{tabular}{ccccc}
\toprule
\texttt{current} & \texttt{temp} & 之後的 \texttt{prev} & 之後的 \texttt{next} & 下一個 \texttt{current} \\
\midrule
1 & \NIL & 2 & \NIL & 2 \\
2 & 1 & 3 & 1 & 3 \\
3 & 2 & 4 & 2 & 4 \\
4 & 3 & 5 & 3 & 5 \\
5 & 4 & 6 & 4 & 6 \\
6 & 5 & \NIL & 5 & \NIL \\
\bottomrule
\end{tabular}
\end{center}
迴圈結束時 \texttt{temp} 是節點 $5$，\texttt{temp->prev} 是 $6$，新的開頭是 $6$。
從 $6$ 沿著 \texttt{next} 走：$6, 5, 4, 3, 2, 1$；沿著 \texttt{prev} 從 $1$ 往回走也
是同樣的順序，兩個方向都接對了。

逐項檢查：
\begin{itemize}
\item (A) 錯。 \texttt{fun} 會把任何串列反轉，$2, 1, 4, 3, 6, 5$ 反轉是
$5, 6, 3, 4, 1, 2$，不是原樣。
\item (B) 錯。 $5, 4, 3, 2, 6, 1$ 反轉是 $1, 6, 2, 3, 4, 5$。
\item (C) 對，上面的追蹤。
\item (D) 錯。 反轉後 $1$ 在最後，不會是 $\ldots 1, 2$。
\end{itemize}

邊界情況：只有一個節點時，迴圈跑一輪，\texttt{temp} 是 \texttt{NULL}，開頭不變，正
確；空串列時迴圈不跑，也正確。 時間 $\Theta(n)$，額外空間 $\Theta(1)$。 和
第~\ref{sec:ll-ops}~節單向串列的 \proc{List-Reverse} 比較：雙向串列每個節點本來就
有兩個方向的指標，對調就好，不用另外記「已經反轉好的一段」。

\pastsol{px:ll-4}{臺科大資工所 112 學年度 第 2 題}
答案：最壞情況 $\Theta(m + n)$（寫成 $O(m + n)$ 也可以）。

\textbf{上界。} 迴圈每一輪做一次（或兩次）比較，每一輪至少讓 \texttt{ptr1}、
\texttt{ptr2} 其中一個前進一格。 \texttt{ptr1} 最多前進 $m$ 次，\texttt{ptr2} 最多
$n$ 次。 設只動 \texttt{ptr1} 的有 $a$ 輪、只動 \texttt{ptr2} 的有 $b$ 輪、兩個都動
的（值相等）有 $c$ 輪，則 $a + c \le m$、$b + c \le n$，輪數是 $a + b + c$。 迴圈在
其中一個變成 \texttt{NULL} 時停：
\begin{itemize}
\item 若最後一輪是值相等、兩個同時變成 \texttt{NULL}，則 $c \ge 1$，輪數
$a + b + c \le (m - c) + (n - c) + c = m + n - c \le m + n - 1$。
\item 否則停下時只有一個走完，例如 \texttt{ptr1} 走完、\texttt{ptr2} 還沒，
$b + c \le n - 1$，輪數 $\le m + n - 1$。 反過來一樣。
\end{itemize}
所以最多 $m + n - 1$ 輪，每輪 $\Theta(1)$，時間 $O(m + n)$。

\textbf{達得到。} 讓兩條串列沒有相同的值，而且全部裡最大和第二大的值分在兩條，例如
list1 是 $1, 2, \ldots, m - 1, m + n$，list2 是 $m, m + 1, \ldots, m + n - 1$。
\texttt{ptr1} 先走過 $1, \ldots, m - 1$（$m - 1$ 輪），停在 $m + n$；之後
\texttt{ptr2} 走過 list2 全部 $n$ 個才變成 \texttt{NULL}。 共 $m + n - 1$ 輪，所以最
壞情況是 $\Theta(m + n)$。 （這和觀念題第 6 題讓 \proc{Merge-Lists} 比較 $m + n - 1$
次的輸入是同一組。）

函式算的是兩條串列共同的值有幾個（兩條各自沒有重複值時，就是交集的大小），走法和
第~\ref{sec:ll-ops}~節的 \proc{Merge-Lists} 一樣，只是不接節點。 最好的情況是一條
串列的值全部比另一條小，$\min(m, n)$ 輪就停，$\Theta(\min(m, n))$。

\pastsol{px:ll-5}{政大資科所 108 學年度 第 9 題}
約定：輸入是串列第一個節點的指標 $h$，每個節點有 $\mathit{next}$；輸出 \TRUE（有環）
或 \FALSE。 用 Floyd 的快慢指標（第~\ref{sec:ll-cycle}~節）：

\begin{pseudo}{政大 108 第 9 題的解答}
\Function{Has-Cycle}{$h$}
    \State $s \gets h$；$f \gets h$ \Comment{慢指標、快指標}
    \While{$f \ne \NIL$ \AND $f.\mathit{next} \ne \NIL$}
        \State $s \gets s.\mathit{next}$ \Comment{走一步}
        \State $f \gets f.\mathit{next}.\mathit{next}$ \Comment{走兩步}
        \If{$s = f$}
            \State \Return \TRUE
        \EndIf
    \EndWhile
    \State \Return \FALSE
\EndFunction
\end{pseudo}

\textbf{正確性。} 做完 $k$ 輪後，$s$ 是從 $h$ 走 $k$ 步到的節點，$f$ 是走 $2k$ 步到
的節點。
\begin{itemize}
\item 沒有環：串列有 $n$ 個節點，$f$ 一直在 $s$ 前面（$2k > k$），兩者是不同的節點，
第 7 行不會執行；$f$ 每輪前進兩格，最多 $\lfloor n/2 \rfloor$ 輪就使 $f$ 或
$f.\mathit{next}$ 變成 \NIL，迴圈停下，回傳 \FALSE。
\item 有環：串列上沒有 \NIL，迴圈只會從第 7 行離開。 設尾巴長 $\mu$、環長 $\lambda$。
第 $\mu$ 輪 $s$ 走到環上，這時 $f$ 早已在環上，沿著環的方向還要追 $g$ 步才碰到 $s$
（$0 \le g < \lambda$）。 之後每輪 $f$ 比 $s$ 多走一步，$g$ 恰好減 $1$，不會跳過
$0$，所以再 $g$ 輪內相遇（$\mu = 0$ 時第 $\lambda$ 輪相遇），回傳 \TRUE。 完整的證
明是定理~\ref{thm:ll-floyd}。
\end{itemize}

\textbf{複雜度。} 沒有環時跑 $\lfloor n/2 \rfloor$ 輪；有環時跑的輪數不超過
$\mu + \lambda = n$。 每輪 $\Theta(1)$，時間 $O(n)$；只用兩個指標，額外空間
$\Theta(1)$，也不改動串列。 題目圖上的串列是整個環（$\mu = 0$、$\lambda = 5$），第
$5$ 輪兩個指標在第一個節點相遇。

另一種寫法是用雜湊表（第~12~章）記錄走過的節點，走到記錄過的節點就有環，走到 \NIL
就沒有；時間一樣是期望 $\Theta(n)$，但要 $\Theta(n)$ 的空間。 題目沒有限制空間時這
也是正確答案，寫 Floyd 的方法比較好。 如果題目還要求環的起點，接著做
\proc{Cycle-Start}（虛擬碼~\ref{alg:ll-floyd}）。
